% paper_18_exchange_constants.tex — GeoVac Paper 18
% Spectral-Geometric Exchange Constants: Projecting Discrete
% Spectra onto Continuous Manifolds
% Status: Draft, April 17, 2026
% Category: Group 3 — Foundations (papers/group3_foundations/)
\documentclass[aps,pra,twocolumn,superscriptaddress,longbibliography]{revtex4-2}

\usepackage{amsmath,amssymb,amsthm,graphicx,xcolor,bm,braket,dcolumn,booktabs}
\usepackage{hyperref}

\newtheorem{theorem}{Theorem}
\newtheorem{observation}{Observation}
\newtheorem{conjecture}{Conjecture}

\begin{document}

\title{Spectral-Geometric Exchange Constants:\\
Projecting Discrete Spectra onto Continuous Manifolds}
\thanks{Paper 18 in the Geometric Vacuum series}

\author{J.~Loutey}
\affiliation{Independent Researcher, Kent, Washington}

\date{April 16, 2026}

%% ========================================================================
%% ABSTRACT
%% ========================================================================

\begin{abstract}
Compact Lie groups have discrete unitary representations
(Peter--Weyl theorem); the GeoVac graph exploits this discreteness
directly.  Transcendental conversion factors appear when these
discrete representations are projected onto non-compact coordinate
systems used for measurement.  We classify these factors by the
type of compactness mismatch they mediate---a six-tier taxonomy
(intrinsic, conformal/calibration, embedding, algebraic-implicit,
composition, inner-factor input data) spanning the GeoVac hierarchy
through Level~5 plus the almost-commutative inner factor---and identify
them with Weyl--Selberg exchange constants of spectral geometry.  The
falsifiable claim: the transcendental content of any GeoVac
observable is determined entirely by the type of projection
linking the compact discrete description to the non-compact
coordinate system.

Beyond the basis-level catalog, three theorems classify the
spectral content of operators on round~$S^3$:
\textbf{(Theorem~1)} operator order is the primary discriminant
of transcendental class---the squared Dirac operator $D^2$ produces
$\pi^{\mathrm{even}}$ exclusively at every integer~$s$, while the
first-order Dirac operator~$D$ produces odd-zeta content
$\zeta_R(2k{+}1) \perp \mathbb{Q}(\pi)$ at odd~$s$;
\textbf{(Theorem~2)} the full classification is three-axis,
adding bundle type (scalar vs.\ spinor, modulating coefficients
only) and diagram topology (free vs.\ vertex-restricted, introducing
Dirichlet $L$-values $G = \beta(2)$, $\beta(4)$ through the
quarter-integer Hurwitz shift);
\textbf{(Theorem~3)} the three components of the
fine-structure-constant combination $K = \pi(B + F - \Delta)$ of
Paper~2~\cite{loutey_paper2} have structurally independent
spectral homes (finite Casimir trace, infinite Dirichlet series at
the packing exponent, Dirac mode degeneracy on the spinor bundle),
each in a distinct cell of the three-axis grid.  The combination
rule $K = \alpha^{-1}$ remains a numerical observation
without first-principles derivation.

A complementary \textbf{Observation (structural incommensurability)}
catalogues the meta-pattern: GeoVac observables that look unitary
at the surface ($K$ itself; two-loop QED on~$S^3$; zero structure
of the GeoVac zeta) resolve under structural analysis into objects
from different cells of the three-axis grid.  No further unification
is currently visible.  The taxonomy is identified as a classification
of Mellin transforms of heat kernels: the free heat kernel produces
$\pi^{\mathrm{even}}$, the interaction-deformed kernel produces
$L$-content via vertex parity, and operator order determines whether
the Mellin integrand samples squared or raw eigenvalues.
The classification is $R$-independent, persisting in the flat-space
limit via Weyl's law.
\end{abstract}

\maketitle

%% ========================================================================
%% I. INTRODUCTION
%% ========================================================================

\section{Introduction}
\label{sec:intro}

The Geometric Vacuum program constructs discrete lattice
Hamiltonians by identifying the natural geometry of each quantum
system and discretizing it as a sparse
graph~\cite{loutey_paper0,loutey_paper1,loutey_paper7}.  A
persistent observation across the hierarchy of natural geometries
is that the discrete algebraic structures---integer eigenvalues,
finite recurrence relations, Gaunt integrals built from Wigner
$3j$ symbols---are entirely self-contained, but that their
connection to physical observables in continuous coordinate
systems always introduces transcendental constants.

The thesis of this paper is that these transcendental conversion
factors are consequences of \emph{compactness}.  Compact Lie groups
have discrete unitary representations (the Peter--Weyl
theorem~\cite{peter_weyl1927}); the GeoVac graph exploits this
discreteness directly.  Transcendental numbers enter when---and only
when---those discrete representations are translated into non-compact
coordinate systems used for measurement.  The taxonomy developed below
classifies the resulting conversion factors by the type of
compact-to-non-compact projection that produces them.  This yields a
falsifiable structural claim: the transcendental content of any GeoVac
observable is determined entirely by the type of projection linking
the compact discrete description to the non-compact coordinate system
(Claim~4, Sec.~\ref{sec:observable_classification}).

The simplest instance is already visible in Paper~0's packing
construction~\cite{loutey_paper0}.  Axiom~2 (geometric indifference)
forces states onto isotropic shells---circles, i.e.\ copies of~$S^1$.
The shell is compact; its circumference $2\pi r$ introduces $\pi$
into the area formula $\sigma_0 = \pi d_0^2/2$ and hence into the
state count $N_k = A_k / \sigma_0 = 2(2k-1)$.  The integer state count
is a consequence of compact topology; $\pi$ is the price of expressing
that count as a ratio of continuous areas.  This pattern---compact
geometry producing integer quantum numbers, with $\pi$ entering
through projection onto non-compact coordinates---recurs at every level
of the hierarchy.

The conversion factors are instances of Weyl's law~\cite{weyl1911} and
the Selberg trace formula~\cite{selberg1956}; what is new is their
systematic identification across the GeoVac hierarchy, culminating in a
falsifiable observable classification (Claim~4) and the observation that
operator order on the compact space determines the motivic weight of the
resulting transcendental content.


%% ========================================================================
%% II. CATALOG OF EXCHANGE CONSTANTS
%% ========================================================================

\section{Catalog of Exchange Constants}
\label{sec:catalog}

\subsection{Level 1: $\kappa = -1/16$ ($S^3 \to \mathbb{R}^3$)}
\label{sec:kappa}

The compactness of $S^3$ is what makes the graph eigenvalues discrete
integers; $\kappa$ converts this discrete spectrum to the non-compact
$\mathbb{R}^3$ energy scale.

The continuum $S^3$ Laplace--Beltrami operator --- which the
(positive-semidefinite) graph Laplacian converges to --- has pure
integer eigenvalues
\begin{equation}
  \lambda_n = -(n^2 - 1), \qquad n = 1, 2, 3, \ldots
  \label{eq:s3_eigenvalues}
\end{equation}
with degeneracy $g_n = n^2$~\cite{loutey_paper7}.  The physical
hydrogen energy levels $E_n = -1/(2n^2)$ are related by the
energy-shell constraint $p_0^2 = -2E$, which acts as a
stereographic focal length projecting $S^3$ onto flat
$\mathbb{R}^3$.  The universal kinetic scale factor
\begin{equation}
  \kappa = -\frac{1}{16}
  \label{eq:kappa}
\end{equation}
calibrates the ground state energy:
$\kappa \cdot \lambda_{\max} = E_H = -1/2$~Ha, where
$\lambda_{\max} \to 8$ is the maximum graph Laplacian eigenvalue
in the continuum limit~\cite{loutey_paper0}.  Per-shell energies
$E_n = -Z^2/(2n^2)$ emerge from the full graph Hamiltonian
$H = \kappa \mathcal{L} + W$ (with node weights $W_{nn} = -Z/n^2$),
not from $\kappa$ alone.

The origin of $\kappa$ is well understood geometrically: it is
the Jacobian of the stereographic projection composed with the
conformal factor of the Fock map.  Specifically, if $\Omega$
is the conformal factor $\Omega = 2p_0/(p^2 + p_0^2)$ (so that
$\Omega(0) = 2$ at the ground-state shell $p_0 = 1$, matching
\S\ref{sec:kappa_derivation} and Paper~7~\cite{paper7}), then
the Laplace--Beltrami operator on $S^3$ and the flat-space
kinetic operator are related by
$\nabla^2_{S^3} = \Omega^{-2}[\nabla^2_{\mathrm{flat}} + \text{lower order}]$,
and $\kappa$ absorbs the accumulated conformal rescaling and
the $p_0$ energy-shell normalization~\cite{loutey_paper7}.

The key property of $\kappa$ is that it is \textit{unnecessary
on the graph}.  The graph Laplacian eigenvalue problem $L\psi =
\lambda\psi$ is solved entirely in terms of integers.  The
constant $\kappa$ appears only when the result is translated
into atomic units (Hartrees, bohrs)---i.e., when the discrete
spectrum is projected onto a continuous coordinate system.

\subsection{Level 2: $e^a E_1(a)$ (Laguerre basis $\to$ prolate spheroid)}
\label{sec:stieltjes}

The Laguerre spectral basis inherits the compact angular structure of
$S^3$; the Stieltjes seed $e^a E_1(a)$ enters when this basis is
embedded in the non-compact prolate spheroidal coordinate, whose
curvature introduces a pole absent from the compact description.

The prolate spheroidal radial solver uses a Laguerre spectral
basis $\varphi_n(\xi) = e^{-\alpha(\xi-1)} L_n(2\alpha(\xi-1))$
with matrix elements computed from three-term
recurrences~\cite{loutey_paper11}.  For $\sigma$ states ($m = 0$),
all matrix elements are purely algebraic: the Laguerre moment
matrices $M_k[i,j] = \int_0^\infty x^k L_i(x) L_j(x) e^{-x}
dx$ are band-structured matrices with rational entries, built
from the recurrence $x L_n = -(n+1)L_{n+1} + (2n+1)L_n -
nL_{n-1}$.

For $m \neq 0$ states, the centrifugal term $m^2/(\xi^2 - 1)$
introduces a $1/x$ singularity in the Laguerre variable
$x = 2\alpha(\xi - 1)$.  The associated Laguerre basis
$L_n^{|m|}(x)$ with weight $x^{|m|} e^{-x}$ absorbs this
singularity by construction.  The centrifugal matrix elements
decompose via partial-fraction expansion:
\begin{equation}
  \frac{m^2}{\xi^2 - 1}
  = 4\alpha^2 m^2 \left[\frac{1}{x}
    - \frac{1}{x + 4\beta}\right],
  \label{eq:partial_fraction}
\end{equation}
where $\beta = \alpha(\xi_0 - 1)$ is the Laguerre scale offset.
The two terms have fundamentally different algebraic character:
\begin{itemize}
  \item The $1/x$ term produces a lowered moment
    $M_{-1}[i,j] = \int_0^\infty x^{|m|-1} L_i^{|m|}(x)
    L_j^{|m|}(x)\,e^{-x}\,dx$, which is fully algebraic via
    the DLMF 18.9.13 summation
    identity~\cite{dlmf_laguerre}.
  \item The $1/(x + 4\beta)$ term produces a Stieltjes matrix
    $J[i,j] = \int_0^\infty \frac{x^{|m|}}{x + a}\,
    L_i^{|m|}(x) L_j^{|m|}(x)\,e^{-x}\,dx$, where $a = 4\beta$.
    This matrix satisfies a three-term recurrence in the
    row/column indices, seeded by the single transcendental
\end{itemize}
\begin{equation}
  S_0^{(\alpha)}(a) = \Gamma(\alpha+1)\,e^a\,E_{\alpha+1}(a),
  \label{eq:stieltjes_seed}
\end{equation}
where $E_n(a) = \int_1^\infty t^{-n} e^{-at}\,dt$ is the
generalized exponential integral and $a = 4\beta$ is the
Laguerre scale parameter~\cite{loutey_track_j}.  For integer
$\alpha = |m|$, all higher $E_n$ reduce to $E_1(a)$ via a
finite rational recurrence, so the entire matrix is built from
algebraic operations plus one transcendental constant.

The computational cost of the transcendental seed is negligible:
one call to $E_1(a)$, which achieves machine precision for
$a \ge 4$ (covering all physical parameters).  The seed then
propagates through an $O(N^2)$ algebraic recurrence that produces
all matrix elements.  The continuous mathematics has been
compressed from a \textit{function} (the quadrature integrand
over $[0,\infty)$) to a \textit{constant} (the Stieltjes seed),
with the entire function-level information reconstructed
algebraically from that constant.

An unexpected computational bonus: the associated Laguerre basis
converges \textit{faster} than the ordinary Laguerre basis for
$m \neq 0$ states.  For $|m| = 1$ ($\pi$ states), the associated
basis is stable by $N = 10$, compared to non-monotonic convergence
at $N \approx 20$ for the ordinary basis.  For $|m| = 2$
($\delta$ states), algebraic-quadrature agreement reaches
$\sim 5 \times 10^{-9}$ (machine precision), with PES shape
(equilibrium geometry) preserved exactly~\cite{loutey_track_j}.
This suggests that identifying the exchange constant and absorbing
it into the basis does not merely explain the transcendence---it
\textit{improves the spectral method}, because the adapted basis
eliminates the curvature mismatch that the exchange constant
mediates.

The Stieltjes seed $e^a E_1(a)$ is the Laplace transform of the
rational function $1/(1+t)$ evaluated at $a$.  It encodes how
the Laguerre inner product (defined on the flat half-line
$[0, \infty)$ with exponential weight) ``sees'' a shifted pole
at $x = -a$ arising from the curvature of the prolate spheroidal
coordinate system.  On the graph, there is no pole; the
centrifugal barrier is a selection rule encoded in the quantum
numbers.  The transcendental appears only when this selection
rule is represented as a rational function in a continuous
coordinate.

\subsection{Level 3: $\mu(R)$ ($S^3 \times S^3 \to$ hyperspherical)}
\label{sec:mu}

On the compact product space $S^3 \times S^3$, the two-electron
Hamiltonian has purely algebraic matrix elements; the adiabatic
eigenvalue $\mu(R)$ enters only when this compact description is
reparameterized into the non-compact hyperspherical coordinate~$R$.

The two-electron problem on $S^3 \times S^3$ is fully algebraic:
the Hamiltonian matrix has entries given by Gaunt integrals
(finite sums of Wigner $3j$ symbols with rational coefficients),
and the eigenvalues are roots of a polynomial with algebraic
coefficients.  The Full CI solver exploits this, achieving
0.35\% error for helium with no transcendentals
anywhere~\cite{loutey_fci}.

The hyperspherical reparameterization introduces coordinates
$(R, \alpha)$ where $R = \sqrt{r_1^2 + r_2^2}$ is the
hyperradius and $\alpha = \arctan(r_2/r_1)$ is the hyperangle.
The angular eigenvalues $\mu(R)$ depend parametrically on $R$
because the effective angular Hamiltonian is
$H(R) = H_0 + R \cdot V^{\mathrm{coupling}}$~\cite{loutey_paper13}.

Each perturbation coefficient $a_k$ in the expansion
$\mu(R) = \mu^{(0)} + a_1 R + a_2 R^2 + \cdots$ has a closed
form: the $a_k$ are finite sums of algebraic numbers (square
roots of rationals from Clebsch--Gordan coefficients) times
powers of $\pi^{-1}$ (from normalization over the angular domain
$[0, \pi/2]$), arising from partial harmonic sums over
$\mathrm{SO}(6)$ Casimir eigenstates~\cite{loutey_paper13}.
For example, $a_1(n{=}1, Z) = (8/3\pi)(\sqrt{2} - 4Z)$.
However, the full function $\mu(R)$ is not transcendental in the
sense originally claimed.  The angular Hamiltonian
$H(R) = H_0 + R \cdot V_C$ is a linear matrix pencil, so its
eigenvalues satisfy the characteristic polynomial
\begin{equation}
  P(R, \mu) \;=\; \det\!\bigl(H_0 + R\,V_C - \mu\,I\bigr) = 0,
  \label{eq:char_poly}
\end{equation}
which is a polynomial of degree $d$ in both $R$ and $\mu$ (where
$d$ is the matrix dimension at a given truncation).  The coefficients
of $P$ lie in $\mathbb{Q}(\pi, \sqrt{2})$, since $H_0$ has integer
entries (SO(6) Casimir eigenvalues) and $V_C$ has entries involving
$\pi^{-1}$ and $\sqrt{2}/\pi$ from the angular integrals.  This
makes $\mu(R)$ an \emph{algebraic} function of $R$ over the
coefficient ring $\mathbb{Q}(\pi, \sqrt{2})$---it is a root of
a polynomial, not a transcendental
function~\cite{loutey_track_p1}.

The transcendental content of $\mu(R)$ is therefore
\emph{inherited} from the coupling matrix $V_C$, which contains
upstream exchange constants of embedding type ($\pi^{-1}$ from
angular normalization, $\sqrt{2}$ from Gaunt-type integrals).
The $R$-dependence itself introduces no new transcendentals:
it is purely algebraic.

The failure of the Taylor/Pad\'e expansion at large $R$
(Track~G~\cite{loutey_track_g}) is reinterpreted as a
\emph{branch point} phenomenon, not evidence of transcendence.
The algebraic curve $P(R, \mu) = 0$ has branch points where
eigenvalue sheets meet---these are the zeros of the discriminant
$\Delta(R) = \mathrm{disc}_\mu(P)$.  The Taylor series around
$R = 0$ has a convergence radius set by the nearest branch point
in the complex $R$-plane.  Pad\'e approximants, being rational
functions, also cannot cross branch points where $\mu(R)$ has
square-root type singularities.  But the implicit relation
$P(R, \mu) = 0$ is exact for \emph{all} $R$, including $R > 5$
where the series diverges.  This has been verified computationally
to machine precision at $R = 50$~bohr~\cite{loutey_track_p1}.

The $O(R) \to O(R^2)$ crossover, which was previously interpreted
as a regime change requiring transcendental interpolation, is
simply the behavior of an algebraic function whose dominant term
changes character: at small $R$, the linear term $V_C \cdot R$
dominates; at large $R$, higher-order terms in the characteristic
polynomial dominate.  The crossover is smooth along the algebraic
curve, even though it is singular from the perspective of a Taylor
expansion centered at $R = 0$.

The crucial point remains that $\mu(R)$ is not intrinsic to the
two-electron problem.  It is an artifact of the hyperspherical
coordinate system.  The $S^3 \times S^3$ product space does not
require $\mu(R)$; it only appears when one reparameterizes for
computational efficiency.  The exchange rate between the discrete
product-space description and the continuous hyperspherical
description is an algebraic function of $R$, with transcendental
content inherited from the angular coupling integrals.


\subsection{Level 4: $\mu(\rho,R)$ (product space $\to$ molecule-frame hyperspherical)}
\label{sec:mu_level4}

At Level~4 the compact product basis (prolate spheroidal CI) is fully
algebraic but structurally incomplete; the non-compact mol-frame
hyperspherical reparameterization introduces $\mu(\rho,R)$ as the
irreducible cost of accessing three-body coalescence physics that the
compact description cannot represent.

Level~4 combines the two-center complication of Level~2 with the
two-electron complication of Level~3.  Two approaches exist, with
fundamentally different algebraic character.

The product-space approach (prolate spheroidal CI, Paper~12) is
fully algebraic: all electron-electron matrix elements are computed
from Neumann expansion recurrence relations with zero spatial
quadrature.  The Laguerre radial and Legendre angular matrix
elements are likewise algebraic.  No transcendental constant
appears anywhere in the computation.

However, a systematic convergence study (Track~M) demonstrates
that this algebraic approach saturates at $\sim$92.5\% of the
exact dissociation energy $D_e$~\cite{loutey_track_m}.  The
angular basis convergence is rapid---92.3\% at $l_{\max} = 2$
(46 basis functions), 92.5\% at $l_{\max} = 4$ (114 basis
functions)---but successive improvements decay by a factor of
200 from $l = 1 \to 2$ to $l = 2 \to 3$, with the $l = 3 \to 4$
increment contributing only 0.04\%.  Extrapolation in the slow
regime predicts that 5\% error would require $l_{\max} \sim 45$
and 1\% error $l_{\max} \sim 113$, both impractical.  The
effective ceiling is 92.5\% $D_e$.

The smoking gun is explicit correlation: 9~basis functions with
$r_{12}^2$ factors (Hylleraas $p = 2$) achieve 94.7\% $D_e$,
exceeding the 114-function prolate CI plateau.  The residual
7.5\% matches Paper~12's cusp diagnosis (7.6\%): the
electron-electron cusp ($1/r_{12}$ singularity) is a completeness
deficiency of the product basis, not a convergence issue.
The transcorrelated approach (Paper~14, Sec.~IV; Paper~15,
Sec.~VI.J) removes this singularity via a Jastrow similarity
transformation, converting $1/r_{12}$ from a hard embedding
constant to a smooth gradient operator at the cost of
non-Hermiticity (Sec.~\ref{sec:taxonomy}).

The mol-frame hyperspherical reparameterization (Paper~15)
resolves this by introducing coordinates $(\rho, R, \alpha)$
where $\rho$ is the electron hyperradius and $\alpha$ the
hyperangle.  The angular eigenvalues $\mu(\rho, R)$ are now a
transcendental function of \textit{two} continuous parameters,
reflecting the simultaneous three-body coalescence in coordinate
and configuration space.  This solver achieves 94.1\% $D_e$,
surpassing the prolate CI ceiling.

The distinction between Levels~3 and 4 is qualitative.  At
Level~3, the $S^3 \times S^3$ FCI converges (albeit slowly)---the
exchange constant $\mu(R)$ is a computational convenience that
accelerates convergence but is not fundamentally necessary.  At
Level~4, the prolate product-space CI \textit{saturates}: the
exchange constant $\mu(\rho, R)$ is irreducible because the
product basis is structurally blind to three-body coalescence.
The reparameterization does not merely accelerate convergence;
it accesses physics that the algebraic product space cannot
represent at any practical basis size.


%% ========================================================================
%% III. THE PATTERN
%% ========================================================================

\subsection{Level 5: Spectral-side exchange constants on $S^3$}
\label{sec:level5_rh_sprint}

Four further realizations of the taxonomy live on the same $S^3$
substrate as Levels 1--4 but classify projections of \emph{spectral}
rather than basis-level data: a closed-form Dirichlet $L$-content
identity in the Dirac Dirichlet difference; a GUE/Poisson dichotomy
between the spectral-zeta and Ihara-zeta sides of the framework; an
algebraic-but-non-radical tier produced by a non-solvable Galois
factor of the Ihara zeta on $S^5$; and an $\alpha^2$-weighted Ihara
zeta living in the ring $\mathbb{Q}(\alpha^2)[s]$.  All four are
catalogued here as new entries in the basis-level catalog of
Levels~1--4; their meta-pattern (categorical disagreement between
spectral and combinatorial views of the same Dirac data) is treated
in Sec.~\ref{sec:meta_observation}.

\paragraph{Closed-form $\chi_{-4}$ content in the Dirac Dirichlet
difference.}  On $S^3$ with the Camporesi--Higuchi Dirac
spectrum $|\lambda_n| = n + 3/2$ and degeneracy $g_n = 2(n+1)(n+2)$,
define the Dirac Dirichlet series $D(s) = \sum_n g_n / |\lambda_n|^s$
and its vertex-parity split $D_{\mathrm{even}}(s)$, $D_{\mathrm{odd}}(s)$
(Paper~28~\cite{loutey_paper28}, eqs.~(3)--(5)). The difference has
the exact closed form
\begin{equation}
  D_{\mathrm{even}}(s) - D_{\mathrm{odd}}(s) \;=\; 2^{s-1}\,\bigl(
    \beta(s) - \beta(s-2)\bigr), \quad s \in \mathbb{Z}_{\ge 2},
  \label{eq:rh_j_identity}
\end{equation}
where $\beta(s) = L(s, \chi_{-4})$ is the Dirichlet $L$-function of
the mod-4 character. Two elementary Hurwitz identities drive the
collapse:
$\zeta(s, 5/4) = \zeta(s, 1/4) - 4^s$ and
$\zeta(s, 3/4) = \zeta(s, 1/4) - 4^s \beta(s)$.
Verified by sympy symbolic identity at $s \in \{2, \ldots, 10, 15\}$
and by PSLQ at 100-digit precision
(\texttt{debug/spectral\_chi\_neg4\_memo.md}).
Eq.~(\ref{eq:rh_j_identity}) is the concrete realization of
Paper~28's qualitative "vertex parity acts as $\chi_{-4}$" statement:
the Dirichlet $L$-content enters the spectral-Dirichlet difference
as a closed-form linear combination of $\beta(s)$ and $\beta(s-2)$
with integer prefactor $2^{s-1}$.

\paragraph{RMT-spacing without critical-line confinement: the
spectral--Ihara dichotomy.}  The zeros of $D(s)$,
$D_{\mathrm{even}}(s)$, $D_{\mathrm{odd}}(s)$ in the strip
$\mathrm{Im}(s) \in [0, 70]$ have unfolded imaginary-part spacing
coefficient of variation $\mathrm{CV} \in \{0.348, 0.343, 0.396\}$
---consistent with GUE ($\mathrm{CV} \approx 0.46$) and inconsistent
with Poisson ($\mathrm{CV} \approx 1.05$). Kolmogorov--Smirnov testing
rejects Poisson at $p = 0.004, 0.030$ respectively and cannot reject
GUE. This is the first GUE signature in the GeoVac framework, and
it is confined to the \emph{spectral} Dirichlet side only: the
combinatorial / Ihara-zeta side (graph non-backtracking walks) has
Poisson / Berry-Tabor spacing at $\mathrm{CV} \approx 1$, a categorical
dichotomy. Equally important, the $D(s)$ zeros are \emph{not} on a
single critical line (they are scattered in
$\mathrm{Re}(s) \in [2.02, 3.42]$), so the GUE signature is a
spacing-only statement, not an RH-style line constraint: classical
$\zeta$ has \emph{both} RMT spacing and critical-line confinement;
GeoVac's spectral zeta has only the first. This creates a new
taxonomic tier, \emph{RMT-spacing without critical-line confinement}.

\paragraph{Algebraic-but-non-radical transcendental content.}
The integer-coefficient Ihara-zeta factorizations produced in
Paper~29~\cite{loutey_paper29} factor over specific number fields as
follows. Cyclotomic factors $\Phi_3$, $\Phi_4$, $\Phi_6$ have Galois
group $\mathbb{Z}/2$ (abelian); they fingerprint specific graph cycles
(e.g.~$\Phi_3, \Phi_6 \leftrightarrow$ 3-prism cycles of the $\ell=1$
Coulomb $S^3$ component). Ramanujan-boundary factors
$(4s^2+1)^2$ and $(9s^2+1)^4$ live in $\mathbb{Q}(i)$, consistent
with the Hopf $\mathbb{Z}_2$ $m$-reflection sub-action of the Paper~25
Hopf $U(1)$ (see Observation~2 of~\cite{loutey_paper29}).
\emph{Bulk} factors have non-abelian Galois groups: $S_3$ for scalar
$S^3$ cubics; $S_4$ for the Dirac-$S^3$ Rule-A $n_{\max}=3$ quartic;
and crucially $S_6$ (non-solvable) for the factor $P_{12}(s^2)$ in
the scalar $S^5$ $N_{\max}=3$ Ihara zeta. By Abel--Ruffini, the 12
zeros of $P_{12}$ have \emph{no radical expression} over
$\mathbb{Q}$; by Kronecker--Weber, they do not lie in any cyclotomic
extension $\mathbb{Q}(\zeta_M)$. This places a new tier in
Sec.~\ref{sec:taxonomy}'s output-class axis: \emph{algebraic but
not radical} --- intermediate between closed-form radical algebraic
and classical periods.

\paragraph{$\alpha^2$-weighted Ihara zeta in the ring
$\mathbb{Q}(\alpha^2)[s]$.}  The unweighted Ihara zeta
depends only on $0/1$ graph connectivity. Introducing the edge
weighting $w(a, b) = 1 + \alpha^2 \cdot |f_{\mathrm{SO}}(a) +
f_{\mathrm{SO}}(b)|/2$, where
$f_{\mathrm{SO}}(n, \kappa) = -(\kappa+1)/[4 n^3 \ell(\ell+1/2)(\ell+1)]$
is the dimensionless Breit--Pauli spin-orbit diagonal, gives a
weighted Ihara zeta $\zeta_G^w(s)$ living in the ring
$\mathbb{Q}(\alpha^2)[s]$ --- an explicit realization of the
spinor-intrinsic tier without the $\gamma = \sqrt{1 - (Z\alpha)^2}$
radical (\texttt{geovac/ihara\_zeta\_weighted.py}; the module's 14 tests
cover the weighted-Ihara infrastructure, not the degree-44 / $-35/648$
values below, which are module-computed but not yet regression-tested ---
a coverage gap logged in \texttt{docs/claim\_test\_matrix.md}).
For Rule-A at $n_{\max}=3$, the inverse zeta expands as a polynomial
in $\alpha$ of degree 44, with $\alpha^2$ coefficient carrying
the clean rational prefactor $-35/648$. The coefficient is a
graph-invariant, not an edge-local perturbation --- a concrete
$\alpha^2$-content projection.

\subsection*{Taxonomic placement of the Level-5 realizations}

The four Level-5 constants slot into the existing taxonomy as
follows. The $\chi_{-4}$ identity (\ref{eq:rh_j_identity}) is a
\emph{flow} exchange constant (infinite spectral sum), with the
novelty that its value is an integer-coefficient linear combination
of classical Dirichlet $L$-values, so Paper~18's transcendental
taxonomy gains a "Dirichlet-$L$-closed" cell at the intersection of
flow and Dirac-spectrum projection. The GUE-vs-Poisson dichotomy is
a \emph{meta-level} discriminant, not itself a transcendental, but
classifies spectral objects by zero statistics --- a new axis in
Sec.~\ref{sec:taxonomy}. The algebraic-but-non-radical tier from
$P_{12}$ is a new \emph{intrinsic} exchange-constant sub-class
between Level~0 (rational/integer) and classical periods; it is the
first time a non-solvable Galois group has been exhibited as a
framework-native transcendental class. The $\alpha^2$-weighted
zeta is a \emph{spinor-intrinsic} exchange constant in the ring
$\mathbb{Q}(\alpha^2)[s]$, a subring of the full spinor-intrinsic
tier $\mathbb{Q}(\alpha^2)[\gamma]/(\gamma^2 + (Z\alpha)^2 - 1)$
introduced in Sec.~\ref{sec:taxonomy}.

% ------------------------------------------------------------------

\section{The Pattern: Weyl--Selberg Exchange}
\label{sec:pattern}

\subsection{Weyl's law}

The mathematical content of the observations in
Sec.~\ref{sec:catalog} is an instance of Weyl's law, which
states that the eigenvalue counting function $N(\lambda)$ of
the Laplace--Beltrami operator on a compact Riemannian manifold
of dimension $d$ satisfies~\cite{weyl1911}
\begin{equation}
  N(\lambda) \sim \frac{\omega_d}{(2\pi)^d}\,
  \mathrm{vol}(\Omega)\,\lambda^{d/2}
  \qquad (\lambda \to \infty),
  \label{eq:weyl}
\end{equation}
where $\omega_d = \pi^{d/2}/\Gamma(d/2+1)$ is the volume of
the unit ball in $\mathbb{R}^d$.

Weyl's law relates two fundamentally different descriptions
of the same manifold:
\begin{itemize}
  \item The \textbf{spectral side}: $N(\lambda)$, a discrete
    staircase function counting integer-labeled eigenvalues.
  \item The \textbf{geometric side}: $\mathrm{vol}(\Omega)$, a
    continuous quantity measured by integration over the manifold.
\end{itemize}
The conversion factor $\omega_d / (2\pi)^d$ involves
$\pi^{d/2}$---a transcendental that depends only on the
dimension.  This factor is the prototypical
\textit{spectral-geometric exchange constant}: it converts
between counting (discrete) and measuring (continuous).

For the specific manifolds in the GeoVac hierarchy, the ball
volumes $\omega_d$ carry $\pi^{\lfloor d/2 \rfloor}$:
\begin{itemize}
  \item $S^1$ ($d = 1$): $\omega_1 = 2$,
    Weyl constant $\omega_1/(2\pi)^1 = 1/\pi$.
  \item $S^2$ ($d = 2$): $\omega_2 = \pi$,
    Weyl constant $\omega_2/(2\pi)^2 = 1/(4\pi)$.
  \item $S^3$ ($d = 3$): $\omega_3 = 4\pi/3$,
    Weyl constant $\omega_3/(2\pi)^3 = 1/(6\pi^2)$.
  \item $S^5$ ($d = 5$): $\omega_5 = 8\pi^2/15$,
    Weyl constant $\omega_5/(2\pi)^5 = 1/(60\pi^3)$.
\end{itemize}
In each case, the Weyl constant involves inverse powers of $\pi$
(specifically $\pi^{-\lceil d/2 \rceil}$), reflecting the
transcendental cost of converting a discrete eigenvalue count
into a continuous volume measure.

\subsection{Compactness as the source of discreteness}

The Peter--Weyl theorem~\cite{peter_weyl1927} states that the unitary
representations of a compact Lie group form a discrete, countable,
orthogonal decomposition of $L^2(G)$.  This is why quantum numbers are
integers: $\mathrm{SO}(4)$ is compact, so its irreps are labeled by
discrete indices $(n, l, m)$ with finite multiplicities $g_n = n^2$.
Exchange constants arise when these discrete representations are
expressed in non-compact coordinate systems---the stereographic
$S^3 \to \mathbb{R}^3$ introduces $\kappa$; embedding in prolate
spheroidal coordinates introduces $e^a E_1(a)$; reparameterizing
$S^3 \times S^3$ as hyperspherical introduces $\mu(R)$.  The
transcendental content is the toll for crossing the compactness
boundary.

\subsection{$\pi$ as the founding example}

The simplest spectral-geometric exchange constant is $\pi$ itself.
The Laplacian on $S^1$ (the circle) has eigenvalues $n^2$, counted
by integers; the continuous circumference is $2\pi$.  The
conversion between the discrete eigenvalue count and the continuous
measure is mediated by $\pi$.  More generally, the volume of the
unit $d$-sphere involves $\pi^{d/2}$, and the Weyl constant
$\omega_d / (2\pi)^d$ carries exactly this factor.  Every exchange
constant in the GeoVac hierarchy is a descendant of this primordial
example: $\kappa$ inherits from the $S^3$ Weyl constant, the
Stieltjes seed involves the exponential integral (a function whose
Taylor coefficients involve the Euler--Mascheroni constant, itself
a limit of harmonic sums weighted by $\pi$-dependent terms), and
$\mu(R)$ is an eigenvalue of a Hamiltonian on $S^5$ whose Weyl
asymptotics carry $\pi^{5/2}$.

\begin{observation}[The graph is $\pi$-free]
\label{obs:pi-free}
Every algebraic operation in the GeoVac graph construction
produces integers or rationals:
\begin{itemize}
  \item The Laplace--Beltrami spectrum $\lambda_n = -(n^2 - 1)$ that the
    graph Laplacian converges to: integer for all $n$.
  \item Degeneracies: $g_n = n^2$, integer.
  \item Gaunt integrals (Legendre normalization,
    $\int P_{l_1} P_k P_{l_2}\,dx$): ratios of factorials via
    Wigner $3j$ symbols, rational.
  \item Selection rules: combinatorial (triangle inequality on angular
    momentum labels).
  \item Clebsch--Gordan coefficients: algebraic numbers (square roots
    of rationals).
  \item Casimir eigenvalues: $l(l+1)$, $\nu(\nu+4)/2$, integer or
    half-integer for all levels.
\end{itemize}
Transcendental numbers ($\pi$, $e^a E_1(a)$, $\zeta(2)$) enter
exclusively through normalization to continuous measures:
$\int |Y|^2\,d\Omega = 1$ (with $d\Omega$ carrying $4\pi$),
$\mathrm{vol}(S^3) = 2\pi^2$, and spectral zeta functions.
The exchange constants are precisely the $\pi$-injection points.
\end{observation}

\subsection{The Selberg trace formula}

Weyl's law is asymptotic.  The Selberg trace
formula~\cite{selberg1956} provides the exact identity:
\begin{equation}
  \sum_j h(r_j)
  = \frac{\mathrm{Area}(\Gamma \backslash \mathcal{H})}{4\pi}
    \int_{-\infty}^{\infty} r\,\tanh(\pi r)\,h(r)\,dr
    + \sum_{\{\gamma\}} \cdots
  \label{eq:selberg}
\end{equation}
where the left side sums over the discrete spectrum, the first
term on the right involves the continuous volume (with a factor
of $\pi$), and subsequent terms sum over conjugacy classes
(geodesics).  The Selberg trace formula is the bridge between
the spectral world and the geometric world, and the
transcendental factors ($\pi$, $\tanh$, logarithms of
algebraic numbers) are the toll for crossing.

When the quotient $\Gamma \backslash G$ is compact, the spectrum
is purely discrete and the trace formula involves only algebraic
data plus $\pi$-dependent conversion factors.  When the quotient
is non-compact, a continuous spectrum appears and the exchange
involves Eisenstein series~\cite{selberg1956}.  This parallels the
GeoVac hierarchy: Level~1 (compact $S^3$) requires only a constant
($\kappa$); Level~3 involves an algebraic function $\mu(R)$ with
transcendental coefficients (Sec.~\ref{sec:algebraic_curve}).

\subsection{The Mellin transform as the taxonomic engine}
\label{sec:mellin}

The preceding subsections identify the \emph{what}---compact-to-non-compact
projection introduces transcendental content---and the \emph{where}---Weyl's
law and the Selberg trace formula are the structural bridges.  Here we
identify the \emph{how}: the Mellin transform of the heat kernel trace is
the single mathematical operation that converts the discrete spectral data
of a compact manifold into the transcendental content classified in
Sec.~\ref{sec:taxonomy}.

For a non-negative self-adjoint operator $A$ on a compact Riemannian
manifold, the spectral zeta function is related to the heat kernel trace
by~\cite{seeley1967,minakshisundaram1949}
\begin{equation}\label{eq:mellin_bridge}
  \zeta_A(s) \;=\; \frac{1}{\Gamma(s)} \int_0^\infty t^{s-1}\,
  \operatorname{Tr}\bigl(e^{-tA}\bigr)\,dt\,,
\end{equation}
where $\operatorname{Tr}(e^{-tA}) = \sum_n g_n\,e^{-t\lambda_n}$ is a
purely spectral object: the sum runs over the discrete eigenvalues
$\lambda_n$ with their integer degeneracies~$g_n$.  Both sides of
Eq.~\eqref{eq:mellin_bridge} encode the same information---the spectrum
of~$A$---but on the left in a form that produces transcendental numbers
($\pi^{d/2}$ from the Seeley--DeWitt asymptotic expansion of the heat
trace, odd-zeta values from the analytic continuation), and on the right
in a form that is manifestly a function of integer-labeled eigenvalues
and degeneracies.  The Mellin transform is the bridge; the $\Gamma(s)$
prefactor is the conversion toll.

\paragraph{Operator order determines the Mellin output class.}
The three-axis classification of QED transcendentals on~$S^3$
(Sec.~\ref{sec:taxonomy}) has its foundation in the distinction between
applying the Mellin transform to the squared eigenvalues $\lambda_n^2$
versus the raw eigenvalues $|\lambda_n|$.

For the \emph{squared} Dirac operator $D^2$ on unit~$S^3$, the
eigenvalues are $(n + 3/2)^2$, which are perfect squares of half-integers.
The heat trace $\operatorname{Tr}(e^{-tD^2})$ is a theta-function-like
object whose Mellin transform samples only even powers: the variable
substitution $m = 2n+3$ converts the spectral sum to a sum over odd
integers, and the Euler--Bernoulli formula for $\zeta_R(2k)$ maps every
term to $\text{rational} \times \pi^{2k}$.  This is the content of the
T9 theorem (Paper~28, Theorem~1~\cite{loutey_paper28}):
$\zeta_{D^2}(s)$ is a two-term polynomial in $\pi^2$ at every
integer~$s$.  No odd-zeta value can appear because the squaring
operation has removed the odd-power structure from the Mellin integrand.

For the \emph{first-order} Dirac operator $|D|$ with eigenvalues
$n + 3/2$, the Mellin transform accesses the raw half-integer spectrum.
The Dirac Dirichlet series $D(s) = \sum g_n |\lambda_n|^{-s}$ involves
both even and odd powers of the spectral parameter, and the parity
discriminant (Paper~28, Theorem~2~\cite{loutey_paper28}) follows
directly: at even~$s$, both constituent Hurwitz zeta terms are
Bernoulli-type ($\pi^{\text{even}}$); at odd~$s$, both involve
genuinely independent odd-zeta values.  The operator order---first
versus second---determines whether the Mellin transform samples the
raw spectrum or its square, and this single distinction generates the
first axis of the three-axis classification.

\paragraph{Interaction deforms the heat kernel.}
The Coulomb potential $V = -Z/r$ perturbs the free Laplacian on~$S^3$.
In the heat kernel framework, the perturbed trace
$\operatorname{Tr}(e^{-t(A + V)})$ has an asymptotic expansion whose
Seeley--DeWitt coefficients $a_k(A + V)$ differ from the free
coefficients $a_k(A)$ by terms involving the potential and its
derivatives~\cite{gilkey1975}.  The Mellin transform of this perturbed
heat trace produces spectral zeta values that access transcendental
content beyond the $\pi$-polynomial ring of the free operator.
(For the \emph{free} Dirac operator~$D$ on unit~$S^3$, the SD
expansion terminates at exactly two terms by the Jacobi theta modular
identity: $\operatorname{Tr} e^{-tD^2}
= (\sqrt{\pi}/2)\,t^{-3/2} - (\sqrt{\pi}/4)\,t^{-1/2}
+ O(e^{-\pi^2/t})$, with $a_k = 0$ for $k \geq 2$; interaction opens
the full MZV algebra.)

The Sommerfeld fine-structure sums $D_p = \sum_{n=1}^\infty c_p(n)$
(Paper~28, Sec.~7~\cite{loutey_paper28}) are precisely the
perturbative expansion of this effect.  At each order $p$ in
$(Z\alpha)^{2p}$, the coefficients $c_p(n)$ involve generalized
harmonic numbers $H_{n-1}^{(r)}$, which are partial Euler sums that
access the full Euler--Zagier multiple zeta value (MZV) algebra.  The
free heat kernel produces only $\pi^{\text{even}}$ (Bernoulli); the
Coulomb interaction opens the door to odd-zeta values ($\zeta_R(3)$
at $p = 2$, $\zeta_R(5)$ and $\zeta_R(7)$ at $p = 3$, and so on)
and to products of zeta values at higher orders.

This maps directly onto the exchange constant taxonomy:
\begin{itemize}
  \item \textbf{Intrinsic tier} (angular/Gaunt structure): no Mellin
    transform is needed; the result is purely rational.
  \item \textbf{Calibration tier} ($\kappa$, $\pi$, $\zeta_R(2)$):
    Mellin transform of the \emph{free} heat kernel on the compact
    manifold $\to$ $\pi^{\text{even}}$ exclusively.
  \item \textbf{Embedding tier} (cusp, $1/r_{12}$, Sommerfeld $D_p$):
    Mellin transform of the \emph{interaction-deformed} heat kernel
    $\to$ the full MZV algebra.
\end{itemize}

\paragraph{The product survival rule as a Mellin selection rule.}
The product survival rule (Paper~28~\cite{loutey_paper28})---the systematic cancellation
of $\zeta(2) \times \zeta(\text{odd})$ products at every Sommerfeld
order---has a natural interpretation in this framework.  The
$n^2$ Fock degeneracy that weights the Sommerfeld sums is the same
$n^2$ that produces $F = \zeta_R(2) = D_{n^2}(4)$ in the free-geometry
Mellin output (Paper~2~\cite{loutey_paper2}, Phase~4F identification).
When the interaction-deformed Mellin transform generates product
terms $\zeta(a) \times \zeta(b)$, the Fock-degeneracy weighting
acts as a selection rule: products involving $\zeta(2)$ are
systematically expelled because $\zeta(2)$ is already ``used'' as
the Fock Dirichlet evaluation $F$, and the $n^2$ weighting enforces
orthogonality between the free and interacting Mellin outputs.

\paragraph{K--Sommerfeld separation as a Mellin consequence.}
The K--Sommerfeld structural separation
(Paper~28~\cite{loutey_paper28})---the PSLQ-verified fact that
$K/\pi$ is not in the $\mathbb{Q}$-span of $\{D_2, \ldots, D_6, 1\}$%
---is a direct consequence of the Mellin framework.
The ingredients of $K = \pi(B + F - \Delta)$ are all
\emph{free-geometry} Mellin invariants: $B$ is a truncated
trace of the free Casimir, $F = \zeta_R(2)$ is the spectral zeta
of the free Fock-degeneracy Dirichlet series at $s = d_{\max}$,
and $\Delta^{-1} = g_3^{\text{Dirac}}$ is a single-level mode
count of the free Dirac operator.  All three live in the
$\pi$-polynomial ring guaranteed by Eq.~\eqref{eq:zeta_d2}.
The Sommerfeld $D_p$, by contrast, are Mellin transforms of the
\emph{Coulomb-perturbed} heat kernel, dressed with harmonic
numbers that access the full MZV algebra.  The two sets of
quantities inhabit different rings because they are Mellin
transforms at different perturbative orders: $K$ at zeroth
order (free geometry), $D_p$ at $p$-th order (interacting geometry).

\paragraph{Summary.}
The Mellin transform \eqref{eq:mellin_bridge} is the silent engine
of the entire exchange constant taxonomy.  It converts discrete
spectral data (integer eigenvalues, integer degeneracies) into
transcendental content ($\pi$, $\zeta(s)$, Dirichlet $L$-values),
and the \emph{class} of transcendental that emerges is determined by
three properties of the operator whose heat kernel is being
transformed: (1)~whether it is first- or second-order (controlling
the motivic weight of the Mellin output); (2)~whether it acts on
scalars or spinors (modulating rational coefficients); and
(3)~whether the potential is free or interaction-deformed (controlling
whether the output stays in the $\pi$-polynomial ring or accesses the
full MZV algebra).  These are precisely the three axes of the QED
transcendental classification formalized in
Sec.~\ref{sec:taxonomy}.

\paragraph{Period-ring stratification of the master Mellin engine
(June 2026 closure).}
The master Mellin engine partition $\mathcal{M}[\mathrm{Tr}(D^k e^{-tD^2})]$
with $k \in \{0, 1, 2\}$ admits a complete period-theoretic
stratification, established by three companion sprints in June~2026
(memos \texttt{debug/sprint\_m1\_pure\_tate\_memo.md},
\texttt{debug/sprint\_mixed\_tate\_test\_memo.md},
\texttt{debug/sprint\_m3\_cyclotomic\_mixed\_tate\_memo.md}):
\begin{equation}
\label{eq:m_engine_period_stratification}
M_1 \subset \mathbb{Q}[\pi, \pi^{-1}], \qquad
M_2 \subset \bigoplus_{k \ge 0} \pi^{2k}\cdot\mathbb{Q}, \qquad
M_3 \subset \mathcal{MT}(\mathbb{Z}[i, 1/2]) \text{ at level } \le 4,
\end{equation}
indexed by the operator-order slot $k$.  M1 (Hopf-base measure
$\mathrm{Vol}(S^2)/4 = \pi$) generates the localised pure-Tate algebra
at $\pi$;\ M2 (Seeley--DeWitt heat-kernel coefficients) inherits the
Fathizadeh--Marcolli~\cite{fathizadeh_marcolli2016} mixed-Tate
classification of spectral-action expansions on Robertson--Walker
spacetimes, with a pure-Tate refinement on the static $S^3$ sub-case
(no $\zeta(3), \zeta(5)$, or MZV content at the SD level);\ M3
(vertex-parity sums producing Hurwitz at quarter-integer shifts and
Dirichlet $L$-values including Catalan $G = \beta(2)$ and $\beta(4)$)
inherits the Deligne~\cite{deligne2010} + Glanois~\cite{glanois2015}
cyclotomic mixed-Tate classification at level $4$, with the
vertex-parity-as-$\chi_{-4}$ correspondence realising the
Deligne--Glanois Galois descent from level $4$ to level $2$ on the
framework's natural QED observables.

Paper~32 §VIII Corollaries \texttt{cor:m1\_pure\_tate},
\texttt{cor:m2\_mixed\_tate}, and \texttt{cor:m3\_cyclotomic\_mixed\_tate}
state the case-exhaustion-theorem-level consequences.  Every
transcendental that has appeared in GeoVac spectral data has a named
motivic home indexed by the Mellin slot~$k$;\ the remainder of this
subsection walks the master-mechanism reading
(\S\,master-mechanism), the mechanism-as-domain sharpening
(\S\,domain-partition), the witness panels for the three period-ring
closures (\S\S\,Mixed-Tate~/~M3-cyclotomic~/~M1-pure-Tate), and the
remaining open questions (sectional convention; non-compact rate;
discrete $c_1$).

\paragraph{Master-mechanism reading (Track~TS-E3, May~2026).}
The Mellin engine reading admits a sharper structural statement.  The
spectral-triple framing of GeoVac (Paper~32) admits a candidate
case-exhaustion theorem (Paper~32 Theorem on $\pi$-source
case-exhaustion): every $\pi$ in a GeoVac observable arises from one
of three named mechanisms M1 (Hopf-base measure), M2 (Seeley--DeWitt
heat-kernel coefficient), or M3 (vertex-topology Hurwitz Dirichlet
$L$-content).  Track~TS-E3's falsification sweep against three
candidate fourth mechanisms (continuum gravity, anomaly classes,
non-trivial principal bundles) returned a stronger structural finding:
M1, M2, M3 are not three independent mechanisms but \emph{three sub-cases
of a single master mechanism} parameterized by operator order:
\begin{equation}
\label{eq:master_mellin}
\pi\text{-source} = \mathcal{M}\!\left[ \mathrm{Tr}(D^k \cdot e^{-tD^2}) \right](s)
\quad\text{with}\quad k \in \{0, 1, 2\}.
\end{equation}
The three sub-cases correspond to:\ $k = 0$ (trivial heat kernel on a
sphere collapses to the Hopf-base measure $\mathrm{Vol}(S^d)$, recovering
M1); $k = 2$ (standard spectral action $\mathrm{Tr}\,f(D^2/\Lambda^2)$
generates Seeley--DeWitt coefficients, recovering M2); and $k = 1$
(Mellin transform of $\mathrm{Tr}(D \cdot e^{-tD^2})$ produces the
$\eta$-invariant series whose vertex-parity decomposition gives the
quarter-integer Hurwitz shifts, recovering M3).  The Mellin engine of
this section is therefore not just \emph{a} mechanism producing
exchange constants; under the case-exhaustion theorem, it is the
\emph{master} $\pi$-injecting mechanism, with operator order $k$
selecting the sub-mechanism.  The first-vs.\ second-order axis named
above (motivic weight) is exactly the $k = 1$ vs.\ $k = 2$ split; the
$k = 0$ free-Mellin case is the Hopf-base measure singled out by Paper~25.

\paragraph{Mechanism-as-domain sharpening (Sprint MR-A/B/C, May~2026).}
The master Mellin engine reading admits a further refinement when read
in the \emph{predictive} direction. The formal statement of the
sharpening appears as a remark accompanying the case-exhaustion theorem
in Paper~32~\S~VIII~\cite{loutey_paper32}; what follows is the
operational summary that closes Paper~18's taxonomic engine. A
three-track sprint sequence (MR-A: Dirac propinquity rate;
MR-B: spectral-action modular residual; MR-C: L2 next-order constant)
tested whether the engine predicts the transcendental ring of a
convergence rate from the operator order $k$ of the underlying mechanism.

Sprint MR-B closed the M2 prediction in closed form. The
Camporesi--Higuchi heat kernel on the half-integer Dirac spectrum
admits a Jacobi $\vartheta_2$ modular decomposition
\begin{equation}
\label{eq:dirac_modular_residual}
   \varepsilon(t) \;=\;
   \sum_{m \ge 1} (-1)^m \sqrt{\pi}\, e^{-m^2 \pi^2 / t}
     \left[ t^{-3/2} - 2 m^2 \pi^2 t^{-5/2} - \tfrac{1}{2} t^{-1/2}\right]
\end{equation}
verified to $> 100$ digits at every tested $t$. Every coefficient sits
in the M2 ring $\sqrt{\pi}\cdot\mathbb{Q} \oplus \pi^2 \cdot \mathbb{Q}$;
the modular exponent is exactly $\pi^2$. The leading prefactor
$2\pi^{5/2}$, the subleading $-\sqrt{\pi}\,t$, and the $\tfrac{1}{2}\sqrt{\pi}\,t^2$
correction at the next order are all in M2.

Sprint MR-A computed the natural Dirac analog of the SU(2) scalar
central Fej{\'e}r propinquity rate. The scalar baseline (L2 sprint,
2026-05-06) had pinned the asymptote rigorously to
$\lim_{n \to \infty} n \gamma_n / \log n = 4/\pi$ with explicit uniform
bound $\gamma_n \le 6\,(\log n)/n$ for $n \ge 2$; the constant
$4/\pi = \mathrm{Vol}(S^2)/\pi^2 = 2\,\mathrm{Vol}(S^1)/\mathrm{Vol}(\mathrm{SU}(2))$
is the Hopf-base measure factor in the natural Haar normalization, the
M1 signature in closed form. The Dirac analog instead obtained the
structural identity $\gamma_n^{\rm Dirac} = \pi$ for every $n_{\max} \ge 1$:
the half-integer-only Peter--Weyl truncation cannot span the constant
function $\chi_0$, so the kernel does not localize at the identity, and
no asymptotic decay exists to PSLQ. The L2 scalar rate decay is in
fact an integer$\times$half-integer interference effect; removing the
integer-spin sector kills it. This is exactly what the M1 mechanism
predicts when the relevant sub-bundle lacks $\chi_0$, and its presence
in the scalar case is precisely what the rigorous $4/\pi$ asymptote
records.

Sprint MR-C extracted the L2 next-order constant $c$ in
$n \cdot \gamma_n^{\rm scalar} \sim (4/\pi)\log(n) + c + O(\log(n)/n)$
to $\geq 80$ digits ($c = 4.1093214674877940927\ldots$), and ran PSLQ
at coefficient ceiling $10^4$ against M1, M2, M3 and the extended
ring $\{\gamma_E, \log 2, G, \zeta(3)\}$. No identification was
returned.

The combined verdict reads:

\begin{quote}
\textit{Mechanism-as-domain partition.}
The index $k \in \{0, 1, 2\}$ classifies both the M-mechanism and
the natural domain of observable from which its signature can be
extracted:
\begin{itemize}
  \item $k = 0$ (M1, Hopf-base measure): state-space propinquity rates;
    signature $4/\pi = \mathrm{Vol}(S^2)/\pi^2$ in the SU(2) Fej{\'e}r
    asymptote (L2 sprint quantitative-rate theorem).
  \item $k = 2$ (M2, Seeley--DeWitt): heat-kernel and spectral-action
    convergence; signature $\sqrt{\pi}\cdot\mathbb{Q} \oplus \pi^2 \cdot \mathbb{Q}$
    throughout the modular residual~\eqref{eq:dirac_modular_residual}.
  \item $k = 1$ (M3, vertex-topology Hurwitz): vertex-restricted
    parity-character spectral sums; signature Dirichlet $\beta(s)$
    content via the closed-form identity
    $D_{\rm even}(s) - D_{\rm odd}(s) = 2^{s-1}(\beta(s) - \beta(s-2))$
    (Paper~28 Theorem~3~\cite{loutey_paper28}); witnessed in
    Paper~28~\S\,QED-vertex (Catalan $G$, $\beta(4)$) but \emph{not} in
    propinquity rates.
\end{itemize}
\end{quote}

Mixing the index of the question (e.g.\ ``propinquity'') with the
index of the mechanism (``vertex parity'') produces structural
degeneracy: this is exactly what MR-A measured. MR-C's negative is
consistent with the same reading: the L2 subleading constant is
downstream of an M1 Stein--Weiss IBP, not a direct
$\mathcal{M}[\mathrm{Tr}(D^k \cdot e^{-tD^2})]$, and the predictive
engine does not commit to the rings of derivative quantities along
subsequent analytical pipelines. The engine is predictive
\emph{within} each domain, where the domain is determined by $k$.
Source memos: \texttt{debug/mr\_\{a,b,c\}\_*\_memo.md}, data
\texttt{debug/data/mr\_\{a,b,c\}\_*.json}.

Sprint Q5'-Stage1-2b (June~2026,
\texttt{debug/sprint\_q5p\_2b\_phi0\_mellin\_memo.md}) sharpens the
$M_1$ sub-mechanism's $\pi$ injection further:\ on a compact
3-dimensional spectral triple, the Hopf-base measure $\pi$ enters
the formal Mellin transform $\Gamma(s) \cdot \zeta_{D^2}(s)$
specifically at the spectral-dimension pole $s = d/2 = 3/2$ via the
Weyl-law short-time leading
$\mathrm{Tr}(e^{-tD^2}) \sim \sqrt{\pi}/2 \cdot t^{-d/2}$;
at every finite truncation $n_{\max}$ the spectral zeta itself is
rational-valued, and the master-Mellin transcendentals appear only
in the continuum analytic-continuation step.  This localizes the
$M_1$ sub-mechanism's natural domain to the non-integer
spectral-dimension residue of the cochain-level Mellin transform.
The same-day continuum follow-on (Sprint Q5'-Stage1-2b-continuum,
\texttt{debug/sprint\_q5p\_continuum\_mellin\_memo.md}) closes the
$M_1$ identification at theorem grade:\ the continuum spectral zeta
$\zeta_{D^2}^{\mathrm{cont}}(s) = 2\,\zeta(2s - 2, 3/2) - \tfrac{1}{2}\,\zeta(2s, 3/2)$
has a simple pole at $s = 3/2$ with bit-exact meromorphic residue $1$,
so $\Gamma(s)\,\zeta_{D^2}^{\mathrm{cont}}(s)$ has Mellin residue
$\Gamma(d/2) = \sqrt{\pi}/2$ at $s = d/2$.  This coefficient is the
same spectral object as the L2 propinquity asymptote
$4/\pi = \mathrm{Vol}(S^2)/\pi^2$ of Paper~38~\S~VIII and as the
Hopf-base Haar normalization $\mathrm{Vol}(S^2)/4 = \pi$ of \S~III.2,
related by exact factors $\pi^{3/2}/8$ and $\pi^2/4$ respectively
with no spectral-data input.  The M2 continuum panel at integer
$s \in \{1, 2, 3, 4, 5\}$ reproduces the Sprint~Q5'-CH-2
Seeley--DeWitt panel bit-exactly.  Tauberian rate-uniformity in a
neighborhood of $s = 3/2$ is the Stage-2-relevant named open gap,
reachable from Karamata~\cite{karamata1962} and
Korevaar~\cite{korevaar2004}.

\paragraph{Connection to motivic / period theory (Sprint Mixed-Tate
Test, June 2026).}
The master Mellin engine reading places the GeoVac taxonomy in
direct contact with the published period / motivic-Galois program
for spectral-action heat-kernel coefficients on continuum
substrates. Fathizadeh--Marcolli~\cite{fathizadeh_marcolli2016}
prove that the coefficients of the asymptotic expansion of the
Connes--Chamseddine spectral action on Robertson--Walker spacetimes
$\mathbb{R} \times S^3$ are periods of mixed Tate motives over
$\mathbb{Q}$. Rejzner's monograph-scale
treatment~\cite{rejzner2016} develops the parallel
period-theoretic framework inside perturbative Algebraic Quantum
Field Theory. Both lines descend from the
Connes--Marcolli~\cite{connes_marcolli2004} renormalization /
motivic-Galois program and the Marcolli--van~Suijlekom gauge-network
spectral-triple framework~\cite{marcolli_vs2014} (with the
Perez-Sanchez correction~\cite{perez_sanchez2024,perez_sanchez2025}
that the continuum limit is Yang--Mills without Higgs), of which
GeoVac is an instance at the data level and an extension at the
structural level.  Brown's mixed-Tate-motives-over-$\mathbb{Z}$
theorem~\cite{brown2012} provides the period-ring backbone
underlying the M3 cyclotomic-mixed-Tate classification. Sprint Mixed-Tate Test (June~2026, memo
\texttt{debug/sprint\_mixed\_tate\_test\_memo.md}) closes the natural
inheritance question for the GeoVac $S^3$ sub-case in the POSITIVE
direction with a pure-Tate refinement:\ at the standard
volume-normalized convention, the GeoVac discrete M2 Seeley--DeWitt
coefficients on $S^3$ sit inside the strictly smaller sub-ring
$\bigoplus_k \pi^{2k}\cdot\mathbb{Q}$ of even-weight pure Tate
periods, with no $\zeta(3)$, $\zeta(5)$, or MZV content appearing
at the SD level. This is sharper than the F--M generic R--W case
in two ways:\ (i) the GeoVac $S^3$ Dirac SD expansion is two-term
exact ($a_k = 0$ for $k \ge 2$, Paper~51 Cor.~2.1), not generically
infinite; (ii) the scalar Laplacian SD coefficients have the
closed-form geometric-series collapse $a_k^\Delta = 2\pi^2/k!$,
not the generic R--W expansion. The M3 sub-mechanism
classification (companion Sprint M3 Cyclotomic Mixed-Tate, June
2026, memo
\texttt{debug/sprint\_m3\_cyclotomic\_mixed\_tate\_memo.md}) places
M3 in $\mathcal{MT}(\mathbb{Z}[i, 1/2])$ at level $\le 4$ at depth
equal to loop order, by direct application of
Deligne~\cite{deligne2010} and the explicit basis of
Glanois~\cite{glanois2015}:\ un-restricted Dirac Dirichlet
$D(s)$ lies in $\mathcal{MT}(\mathbb{Z})$ at depth $1$;
vertex-restricted $D_{\mathrm{even}} - D_{\mathrm{odd}}$ lies in
$\mathcal{MT}(\mathbb{Z}[i, 1/2])$ at level $4$ depth $1$ via
quarter-integer Hurwitz $+$ Dirichlet $\beta$ (Catalan $G$,
$\beta(4)$);\ the depth-$k$ loop tower lives at depth $k$ in
$\mathcal{MT}$ at level $\le 4$.  M3 does \emph{not}
contaminate the SD coefficients (M3 lives at $k=1$ in
$\mathcal{M}[\mathrm{Tr}(D^k e^{-tD^2})]$, M2 at $k=2$).  Together
with the companion Paper~32~\S~VIII corollaries
\texttt{cor:m2\_mixed\_tate} and \texttt{cor:m3\_cyclotomic\_mixed\_tate},
the master Mellin engine now has a complete period-theoretic
stratification:\ M1 in $\mathbb{Q}[\pi, 1/\pi]$, M2 in
$\bigoplus_k \pi^{2k}\cdot\mathbb{Q}$, M3 in
$\mathcal{MT}(\mathbb{Z}[i, 1/2])$ at level $\le 4$.

\paragraph{M3 sectional convention refinement (Sprint L2-D, May 2026).}
The vertex-parity Hurwitz mechanism of M3 is sectional in the chosen
index parity convention, NOT in spacetime signature.  Specifically:
at the Krein extension to $(3, 1)$ (Sprint L2-D, Paper~32
\S~VIII.D~\cite{loutey_paper32}), M3's L0
trivialization prediction holds trivially under chirality-pairing
(chirality symmetry makes $D_{\mathrm{even}} = D_{\mathrm{odd}}$
identically) but falsifies under Paper~28's
$n_{\mathrm{Fock}}$-parity reading --- the convention that accesses
Catalan $G = \beta(2)$ via the closed-form identity
$D_{\mathrm{even}}(s) - D_{\mathrm{odd}}(s) = 2^{s-1}(\beta(s) -
\beta(s-2))$ (Paper~28 Theorem~3~\cite{loutey_paper28}).  The
Bizi--Brouder--Besnard sign-flip $\{J, \gamma^5\} = 0$ at
$(m, n) = (4, 6)$ is a spectral-triple-axiom statement, not a
vertex-parity-sum statement.  M3's transcendental content
($\beta(2)$, $\beta(4)$, Catalan $G$, Dirichlet $L$-values) is
therefore tied to the $n_{\mathrm{Fock}}$-parity sector choice of
the truncation, not to the signature of the underlying spectral
triple.  This refinement does not change the master Mellin engine
domain partition above (M1 $\leftrightarrow$ propinquity rates,
M2 $\leftrightarrow$ heat-kernel and spectral-action convergence,
M3 $\leftrightarrow$ vertex-restricted parity-character spectral
sums); it sharpens the parity-convention-dependence of the M3
sub-mechanism's accessibility.  Source memos:
\texttt{debug/l2\_d\_connes\_axiom\_audit\_31\_memo.md},
data \texttt{debug/data/l2\_d\_connes\_axiom\_audit\_31.json}.

\paragraph{Operator-vs.-slot refinement of the M3 sub-mechanism
(Sprint NA-1, June 2026).}
The Mellin-slot index $k \in \{0, 1, 2\}$ in
$\mathcal{M}[\mathrm{Tr}(D^k\, e^{-tD^2})]$ classifies the
\emph{sub-mechanism} (M1 / M3 / M2); it does \emph{not} by itself fix
the period ring of the output.  The Mellin slot exponent~$s$ --- the
argument at which the moment is evaluated --- selects \emph{which
periods inside the sub-mechanism's natural period ring} are accessed.
Sprint NA-1 (\texttt{debug/sprint\_na1\_depth2\_mellin\_memo.md},
2026-06-06) makes this distinction explicit on the natural
Camporesi--Higuchi diagonal substrate for the $k = 1$ slot (M3).
Let
\begin{equation}
\label{eq:m3_gammaP_slot}
M_3^{\gamma_P}(r)
   \;:=\; \sum_{n \ge 0} (-1)^n\, g_n\, |\lambda_n|^{-r}
\end{equation}
be the parity-graded Mellin function of the Camporesi--Higuchi Dirac
spectrum, with $\gamma_P|n\rangle = (-1)^n |n\rangle$ the vertex-parity
grading on the eigenshells.  The \emph{same operator} $\gamma_P$
produces two structurally distinct period-ring outputs depending on
the Mellin slot exponent at which it is probed:
\begin{itemize}
  \item \textbf{Pure-Tate slot} ($r = 2s - 1 \in 2\mathbb{Z} + 1$, odd
    integer; Paper~55~\cite{loutey_paper55}
    \S\,M3-diagonal-collapse Theorem on
    $M_3^{\gamma_P}$ closed form):
    closed form
    $M_3^{\gamma_P}(s) = 2^{2s - 3}\,(\beta(2s - 1) - \beta(2s - 3))$,
    landing in $\bigoplus_k \pi^{2k+1}\cdot\mathbb{Q}$ via the
    Euler-number closed form
    $\beta(2k + 1) \in \pi^{2k + 1}\cdot\mathbb{Q}$.  This is
    \emph{pure-Tate at level $1$}.
  \item \textbf{Level-$4$ cyclotomic slot} (integer $s$ on the
    un-shifted Dirac Dirichlet difference; Paper~28
    Theorem~3~\cite{loutey_paper28}):
    $D_{\mathrm{even}}(s) - D_{\mathrm{odd}}(s)
       = 2^{s-1}(\beta(s) - \beta(s-2))$,
    accessing $\beta(\mathrm{even})$ values
    (Catalan $G = \beta(2)$, $\beta(4)$) inside
    $\mathcal{MT}(\mathbb{Z}[i, 1/2])$ at level $4$
    (Paper~55~\cite{loutey_paper55} \S\,M3-Galois-descent
    Proposition on vertex parity as
    Deligne--Glanois descent).
\end{itemize}
Both probes act on the same diagonal parity grading $\gamma_P$
sitting in the $k = 1$ slot of the master Mellin engine; the
period content (pure-Tate vs.\ level-$4$ cyclotomic) is selected
by the Mellin slot exponent~$r$, not by a different operator.  The
Glanois~\cite{glanois2015} basis~$\mathcal{B}^4$ accommodates both
columns:\ the pure-Tate column is the $N = 1$ sub-basis, the
level-$4$ column is the genuinely $N = 4$ content.

This refinement does not change the case-exhaustion statement
(Paper~32~\S~VIII), the M3 sectional convention refinement
above, or the period-ring stratification
$M_3 \subset \mathcal{MT}(\mathbb{Z}[i, 1/2])$ at level $\le 4$.
It sharpens the description of the M3 row of the master Mellin
engine:\ the sub-mechanism index $k = 1$ commits to the
$\eta$-invariant Mellin family
$\mathcal{M}[\mathrm{Tr}(D \cdot e^{-tD^2})]$, but the period-ring
location within $M_3$ at each evaluation is determined by the slot
exponent.  Cross-reference Paper~55~\cite{loutey_paper55}
\S\,M3-diagonal-collapse and \S\,M3-Galois-descent for the
explicit pure-Tate vs.\ level-$4$ demonstration, including the
NA-1 diagonal-collapse identity
$J(s_1, s_2) = M_3^{\gamma_P}(s_1 + s_2 - 1)$ on the
Camporesi--Higuchi diagonal substrate.

An identical operator-vs.-slot reading is available for M2 (the
Seeley--DeWitt slot, $k = 2$):\ the squared-Dirac heat-kernel
operator $D^2$ is fixed, but the Mellin slot exponent~$s$ selects
which Seeley--DeWitt coefficient $a_k(D^2)$ is read off via
$\mathrm{Res}_{s = (d - 2k)/2}\,\zeta_{D^2}(s)$.  The pure-Tate
M2 closure
$M_2 \subset \bigoplus_k \pi^{2k}\cdot\mathbb{Q}$
(Sprint Mixed-Tate Test, paragraph above) describes the period
ring;\ which $\pi^{2k}\cdot\mathbb{Q}$ slot is accessed by a given
observable is fixed by the Mellin slot exponent at which the
heat-kernel trace is evaluated.  The M1 slot $k = 0$ is degenerate
in this sense:\ the trivial heat kernel collapses to the
Hopf-base measure $\mathrm{Vol}(S^d)$, and the period ring
$\mathbb{Q}[\pi, \pi^{-1}]$ has no internal Mellin-slot stratification
to select.  The operator-vs.-slot distinction is therefore
non-trivial for M2 and M3 (whose period rings have internal weight
gradings) and degenerate for M1 (whose period ring has only one
graded piece, $\mathbb{Q}[\pi, \pi^{-1}]$ at depth~$0$).

\paragraph{Why slot, not operator (Sprint NA-1-offdiag, June~2026).}
The structural reason the operator-vs.-slot distinction works as
stated above:\ the trace functional
$\mathrm{Tr}(\cdots\, e^{-t_i D^2}\, \gamma\, e^{-t_j D^2}\cdots)$
commits to the \emph{diagonal} of any inserted operator $\gamma$ in
the $D$-eigenbasis and loses the off-diagonal information.  For any
self-adjoint $\gamma$ on a finite-dimensional Hilbert space ---
commuting or non-commuting with $D$ ---
\[
\mathrm{Tr}\big(D^2\, e^{-t_1 D^2}\, \gamma\, D\, e^{-t_2 D^2}\big)
\;=\; \sum_i \widetilde{\gamma}_{ii}\,
\lambda_i^3\, e^{-(t_1 + t_2)\lambda_i^2},
\]
depending only on $t_1 + t_2$ (Paper~55~\cite{loutey_paper55}
\S\,M3-diagonal-collapse, Theorem on the off-diagonal CH substrate).
The substrate change from a diagonal $\gamma$-grading to a
non-diagonal one (e.g.\ chirality-flipping $E_1$ entries on the full
spinor bundle) does not change the trace output, because off-diagonal
entries of $\widetilde\gamma$ are invisible to the trace.  This is
the deeper structural reason the operator-vs.-slot reading commits
to the slot exponent:\ the operator side of the master Mellin engine
is projected onto its diagonal in the spectral basis before any
period content is read off, and only the slot exponent~$s$ then
selects which period ring element is accessed.

\paragraph{Compact / non-compact partition of the $4/\pi$ rate
(Sprint M-Z, May 2026).}
The master Mellin engine's $M_1$ Hopf-base measure mechanism produces
the universal $4/\pi$ rate constant in the propinquity convergence of
Paper~38 (compact $SU(2)$ truncation) and in the universal rate
theorem of Paper~40 (general compact connected Lie groups with
bi-invariant metric).  Sub-sprint M-Z (modular propinquity
investigation, 2026-05-23) flagged a structural reason why this rate
does \emph{not} transport to non-compact Coulomb (Sturmian basis on
$L^2(\mathbb{R}^3)$): the Hopf-base measure mechanism is compact-Lie-quotient
specific; Sturmian's continuum content dilutes the Plancherel weight
to 1, killing the localization that drives the
$\log n / n$ rate on $SU(2)$.  A rigorous non-compact analog of
the Paper~38 / Paper~40 rate remains an open question; Latr\'emoli\`ere
2018~\cite{latremoliere2018} establishes existence of a complete
propinquity metric on the inductive-limit construction but not a rate.
Source: \texttt{debug/sprint\_modular\_propinquity\_mZ\_bethe\_log\_memo.md}.

\paragraph{M1 period-ring closure (Sprint M1 Pure-Tate, June 2026).}
The M1 sub-mechanism produces period values in the localised
pure-Tate sub-ring $\mathbb{Q}[\pi, \pi^{-1}]$ of mixed Tate periods
over $\mathbb{Q}$, at depth $0$ and Tate weight $\in \mathbb{Z}$, by
direct identification of $\pi = \mathrm{Vol}(S^2)/4$ as a weight-$1$
Tate period (Kontsevich--Zagier 2001).  Every M1 witness logged
across the corpus --- $\mathrm{Vol}(S^2)/4 = \pi$ (Paper~25
Hopf bundle base measure); $4/\pi = \mathrm{Vol}(S^2)/\pi^2$
(Paper~38 §L2 asymptotic propinquity rate); $4/\pi$ (Paper~40 Thm~1
universal rate constant across compact Lie groups); $1/\pi^2$
(Paper~43 §10.2 Pythagorean orthogonality);\ $1/\pi^2$ inside
$-3\zeta(3)/(16\pi^2)$ (Paper~50 §3 F-theorem, an M1$\times$M3
cross-product) --- is a $\mathbb{Q}$-linear combination of $\pi^n$
for $n \in \mathbb{Z}$.  Combined with this sprint's two companion
closures (Sprint Mixed-Tate Test placing M2 in
$\bigoplus_k \pi^{2k}\cdot\mathbb{Q}$, and Sprint M3 Cyclotomic
Mixed-Tate placing M3 in $\mathcal{MT}(\mathbb{Z}[i, 1/2])$ at level
$\le 4$), the master Mellin engine of \S~III.7 has a complete
period-theoretic stratification indexed by the Mellin slot
$k \in \{0, 1, 2\}$:
\begin{equation*}
M_1 \subset \mathbb{Q}[\pi, \pi^{-1}], \quad
M_2 \subset \bigoplus_k \pi^{2k}\cdot\mathbb{Q}, \quad
M_3 \subset \mathcal{MT}(\mathbb{Z}[i, 1/2]) \text{ at level} \le 4.
\end{equation*}
Companion Paper~32 §VIII Corollary \texttt{cor:m1\_pure\_tate} states
the case-exhaustion-theorem-level consequence; memo
\texttt{debug/sprint\_m1\_pure\_tate\_memo.md}.

\paragraph{Discrete $c_1$ verification (Sprint TS-E3 falsification target).}
The strongest concrete computational falsification target flagged by Track
TS-E3 is the discrete first Chern number of the Hopf U(1) bundle on the
finite Fock-projected $S^3$ graph: under the master Mellin engine reading,
the discrete $c_1$ should be \emph{integer-valued} on any finite cutoff,
because the construction engages no continuous integration step (no
Hopf-base measure, no heat-kernel trace, no Mellin transform), and
therefore none of M1, M2, M3.  Both physically natural $U(1)$
phase conventions (Condon--Shortley canonical, all phases zero; and Z$_2$
alternating, $\phi = \pi$ on $L_\pm$ edges) yield exactly $c_1 = 0$ at
$n_{\max} = 2$ and $n_{\max} = 3$, confirmed in exact rational arithmetic
(Fraction-typed phases, no floating-point ambiguity).  A sanity-check
convention (prescribed flux $1/4$ per plaquette via non-tree-edge phases)
recovers $|c_1| = (\#\text{cycles}) / 4 = 1/2$ at $n_{\max} = 3$ exactly,
and a stress-test convention ($1/3$ flux) recovers $|c_1| = 2/3$,
confirming the algorithm tracks rational non-integer fluxes faithfully
and would surface a non-integer result if one existed under the natural
conventions.  The continuum first Chern integer reappears via the M1
$\mathrm{Vol}(S^2)/4 = \pi$ normalisation in the continuum limit; on the
finite graph the topological invariant is integer in $\mathbb{Q}$.  Source:
\verb|debug/discrete_c1_hopf_s3.py|, data \verb|debug/data/discrete_c1_hopf_s3.json|,
memo \verb|debug/discrete_c1_falsification_memo.md| (Sprint TS-E3, 2026-05-04).


%% ========================================================================
%% IV. TAXONOMY
%% ========================================================================

\section{Taxonomy of Exchange Constants}
\label{sec:taxonomy}

The following taxonomy classifies exchange constants by the type of
compactness mismatch they mediate.  The hierarchy---intrinsic,
conformal/calibration, embedding, algebraic-implicit, composition,
and inner-factor input data---can be read as increasing departure from
compactness: intrinsic constants convert between compact manifolds of
different dimension; conformal/calibration constants convert between a
compact manifold and flat space; embedding constants arise from
coordinate singularities in the non-compact description;
algebraic-implicit functions encode parameter-dependent geometry
through a characteristic polynomial whose coefficients are inherited
from upstream constants; composition constants measure the cost of
factorizing a single compact space into multiple incompatible
coordinate systems; and inner-factor input data are the calibration
constants of the almost-commutative inner factor (the Standard-Model
Yukawa/Higgs sector), supplied externally rather than generated by any
projection.

Based on the catalog and the Weyl--Selberg identification, we
propose the following taxonomy:

\begin{table}[h]
\caption{Spectral-geometric exchange constants in the GeoVac
  hierarchy.  ``Source'' indicates the discrete structure;
  ``Target'' indicates the continuous coordinate system.
  ``Type'' classifies the exchange constant by what determines it.
  ``Algebraic (implicit)'' indicates that the function satisfies a
  polynomial equation $P(R,\mu) = 0$ with coefficients inherited from
  upstream exchange constants (Sec.~\ref{sec:algebraic_curve}).
  ``Piecewise'' indicates that the angular Hamiltonian has
  piecewise-smooth parameter dependence (from the split-region
  Legendre expansion), preventing a global characteristic polynomial.
  The function is piecewise-algebraic within each region.}
\label{tab:taxonomy}
\begin{ruledtabular}
\begin{tabular}{ccccc}
  Level & Source & Target & Constant & Type \\
  \hline
  --- & $S^1$ lattice & $\mathbb{R}^1$ & $\pi$ & intrinsic \\
  1 & $S^3$ graph & $\mathbb{R}^3$ & $\kappa = -1/16$ & conformal \\
  2 ($m{=}0$) & Laguerre rec. & prolate sph. & (none) & algebraic \\
  2 ($m{\neq}0$) & assoc. Laguerre & prolate sph. & $e^a E_1(a)$
    & embedding \\
  3 (FCI) & $S^3 {\times} S^3$ & (none) & (none) & algebraic \\
  3 (hyper.) & $\mathrm{SO}(6)$ reps & $S^5(R)$ param.
    & $\mu(R)$ & alg.\ (implicit) \\
  4 (prolate CI) & prolate product & (none) & (none) & algebraic \\
  4 (hyper.) & $\mathrm{SO}(6)$ reps & $(\rho,R)$ param.
    & $\mu(\rho,R)$ & piecewise \\
  5 (composed) & Level 3+4 fibers & full $N$-electron
    & PK pseudopot. & composition \\
  5 (balanced) & hydrogenic $\times$ multipole & cross-center $V_{ne}$
    & $\gamma(k,\alpha R)$ & algebraic \\
\end{tabular}
\end{ruledtabular}
\end{table}

The taxonomy reveals a hierarchy of exchange types:
\begin{enumerate}
  \item \textbf{Algebraic} (no exchange constant): When the
    computation stays on the graph or uses a basis whose inner
    product matches the graph structure, no transcendental
    appears.  Examples: Level~1 on $S^3$, Level~2 $\sigma$
    states, Level~3 FCI.

  \item \textbf{Intrinsic constants} ($\pi^{d/2}$): Determined
    by the manifold's dimension and volume alone, independent of
    any discretization or basis choice.  These are the Weyl
    constants---universal volume-to-eigenvalue conversion factors.
    Example: $\pi$ itself, for $S^1 \to \mathbb{R}$.

  \item \textbf{Conformal constants}: Determined by the conformal
    geometry of the projection from the compact graph manifold to
    flat space.  Example:
    $\kappa = -1/16 = -1/\Omega^4(0)$, where $\Omega(0) = 2$ is
    the stereographic conformal factor at $p = 0$
    (Sec.~\ref{sec:kappa_derivation}).  Equivalently derivable
    from graph vertex degree $d_{\max} = 4$; both routes yield
    the same value because the graph discretizes $S^3$.

  \item \textbf{Embedding constants}: Single transcendental
    numbers seeding algebraic recurrences.  Determined by the
    interaction between a spectral basis and a coordinate
    singularity.  Examples: $e^a E_1(a)$ for the prolate
    centrifugal pole; the $1/r_{12}$ electron-electron cusp,
    which is the embedding exchange constant for the two-electron
    coalescence in any product basis.
    The cusp embedding constant is unique in being \emph{removable}
    by a similarity transformation: the transcorrelated (TC)
    Jastrow factor $J = \tfrac{1}{2} r_{12}$ satisfies
    $e^{-J} (1/r_{12})\, e^{J} = \text{smooth}$, converting
    the algebraic partial-wave convergence $1/(l+\tfrac{1}{2})^4$
    to exponential.  This makes $1/r_{12}$ a ``soft'' embedding
    constant---present in any product-space representation but
    eliminable by a non-unitary change of basis---in contrast to
    $e^a E_1(a)$, which is structurally irreducible.

    The $1/r_{12}$ classification is \emph{universal}: the cusp
    condition survives relativistic corrections with only an
    $O((Z\alpha)^2)$ amplitude rescaling (Kutzelnigg-Morgan 1992), and
    partial-wave convergence exponents ($l^{-4}$ singlet, $l^{-6}$
    triplet) are identical in Schr\"odinger-Coulomb and Dirac-Coulomb
    (verified numerically at $Z=4$).  The Dirac program does not
    regularize the e-e cusp; full QED would be required.

    A \emph{distributional} sub-category appears in the Breit spin-spin
    interaction: the bare $1/r_{12}^3$ kernel diverges logarithmically
    at coalescence for every partial wave.  The Breit-Pauli tensor
    projection regularizes this, leaving $K_l^{BP} = r_<^l / r_>^{l+3}$
    with exact rational matrix elements for same-$n$ orbitals
    ($R_{BP}^2(2p,2p;2p,2p) = 7/768$ at $Z=1$).  For cross-$n$
    orbitals a new transcendental appears:
    $M_{\rm ret}^2(1s,2p;1s,2p) = 21344/729 - (10240/243)\log(2)$,
    the first $\log$-embedding constant in the framework.

    On the $S^3$ pair-state lattice (v2.9.1), 92--94\% of $V_{ee}$'s
    Frobenius mass is diagonal in the graph eigenbasis, with the
    off-diagonal residual concentrated on the $(1s,1s)$ coalescence
    pair-state.  The graph nearly diagonalizes $V_{ee}$ without
    containing $1/r_{12}$---direct evidence that the cusp is embedding
    content the graph cannot absorb, explaining the graph-native CI
    convergence floor (Paper~13, Track~DI).

  \item \textbf{Algebraic (implicit) functions}: Functions of
    continuous parameters defined implicitly by $P(R,\mu) = 0$,
    where $P$ is the characteristic polynomial of the
    parameter-dependent Hamiltonian $H(R) = H_0 + R \cdot V_C$.
    The coefficients of $P$ lie in $\mathbb{Q}(\pi, \sqrt{2})$---a
    transcendental extension of $\mathbb{Q}$---but the $R$-dependence
    itself is algebraic: $\mu(R)$ is a root of a polynomial in $R$.
    The transcendental content is \emph{inherited} from the angular
    coupling matrix $V_C$ (which contains upstream exchange constants
    of embedding type), not generated by the parameterization.
    Example: $\mu(R)$ for the $\mathrm{SO}(6) \to \mathrm{SO}(4)$
    symmetry crossover on $S^5$ at Level~3.  At $l_{\max} = 0$,
    $\mu(R)$ is an explicit linear function of $R$; at higher
    $l_{\max}$, it is an algebraic function of degree $l_{\max}+1$,
    with branch points at eigenvalue crossings in the complex
    $R$-plane.  At Level~4, the split-region Legendre expansion of
    the nuclear charge function introduces piecewise-smooth
    $\rho$-dependence at the region boundary $\rho = 1/2$, breaking
    the linear pencil structure. The angular eigenvalues $\mu(\rho)$
    are piecewise-algebraic within each region but do not satisfy a
    single global characteristic polynomial. The product-space CI
    saturation at 92.5\% $D_e$ remains the key irreducibility result.

  \item \textbf{Composition constants}: The accuracy cost of
    factorizing a full $N$-electron wavefunction into composed
    coordinates.  Unlike the preceding types, which convert between
    discrete and continuous descriptions, composition constants
    convert between a single-geometry and a multi-geometry
    (factorized) description.  The factorization destroys inter-group
    antisymmetry---the exchange operator $P_{ij}$ mixes coordinate
    systems that composition separates---and the PK pseudopotential
    is the local approximation to this nonlocal effect.  Example:
    PK for LiH composed geometry (Level~3 core + Level~4 valence),
    which achieves 144$\times$ angular compression at the cost of
    ${\sim}5\%$ $R_{\rm eq}$ error at $l_{\max} = 2$.

  \item \textbf{Inner-factor input data}: The calibration constants of
    the almost-commutative inner factor---parameter-tied Dirichlet
    rings generated by the eigenvalues of a finite Dirac operator
    $D_F$ on a Connes--Chamseddine-compatible inner factor (the
    Standard-Model Yukawa/Higgs sector).  Unlike the preceding types,
    which a projection \emph{generates} from the discrete graph,
    inner-factor data are empirical input from outside the framework:
    GeoVac admits a Higgs sector structurally but does not
    autonomously select the Yukawa couplings (Sprint~H1; Paper~32
    \S~VIII.C), so these constants are inputs, not outputs.  They are
    categorically disjoint from the outer-factor period content (the
    $M_1/M_2/M_3$ Mellin sectors) and tensor cleanly with it, neither
    determining the other.
\end{enumerate}

The dependency hierarchy is explicit: intrinsic constants depend
only on dimension; conformal/calibration constants depend on the
graph; embedding constants depend on the basis; algebraic-implicit
functions depend on the full parameterized geometry; composition
constants depend on the factorization topology (which groups share
coordinates); and inner-factor input data depend on the choice of
finite spectral triple, supplied from outside the framework.

\paragraph{Odd-zeta content from first-order spectral operators.}
The Dirac operator on unit $S^3$ (Camporesi--Higuchi~\cite{camporesi_higuchi1996}
spectrum $|\lambda_n| = n + 3/2$, degeneracy
$g_n^{\text{Dirac}} = 2(n+1)(n+2)$) introduces a structurally new
transcendental: odd-zeta values $\zeta_R(2k{+}1)$ enter the Dirac
Dirichlet series $D(s) = \sum_n g_n^{\text{Dirac}}/|\lambda_n|^s$ at
odd integer $s \ge 5$.  The Hurwitz decomposition (Theorem~1) gives
\begin{equation}\label{eq:dirac_zeta3_witness}
  D(5) \;=\; 14\,\zeta_R(3) \;-\; \tfrac{31}{2}\,\zeta_R(5),
\end{equation}
the first integer~$s$ at which $\zeta_R(3)$ appears with a non-divergent
coefficient (at $s = 3$ the term proportional to $\zeta_R(s{-}2) = \zeta_R(1)$
diverges).  At the packing exponent $s = 4$, by contrast,
\begin{equation}\label{eq:dirac_d4_value}
  D(4) \;=\; 6\,\zeta_R(2) - \tfrac{15}{2}\,\zeta_R(4)
        \;=\; \pi^2 - \frac{\pi^4}{12},
\end{equation}
purely $\pi^{\text{even}}$---consistent with Theorem~\ref{thm:T9}
below (squaring the spectrum maps every odd-power Dirichlet sum to an
even power, and $\zeta_{D^2}(s)$ lies in $\sqrt{\pi}\cdot\mathbb{Q} \oplus \pi^2\cdot\mathbb{Q}$
at integer $s$).  The half-integer eigenvalue $|\lambda_n| = n + 3/2$
is essential: naive substitution of an integer~$m$ for the eigenvalue
would predict an even-$s$ odd-zeta witness, contradicting Theorem~1
part~(2) and Paper~28~\cite{loutey_paper28} Theorem~5 / Eq.~(43).
The shift~$3/2$ is what cancels the singular $m=1$ term in the
odd-integer rewrite of the spectral sum.

This odd-zeta content (entering at odd $s$ on the spinor bundle) is
categorically different from calibration-tier content.  Calibration
constants carry even-zeta or $\pi^{\text{even}}$ content arising from
second-order Riemannian operators through Jacobi-$\theta$ inversion;
the odd-zeta $\zeta_R(3)$ enters through the first-order Dirac operator
via half-integer Hurwitz-zeta identities at \emph{odd} $s$, and is
$\mathbb{Q}$-linearly independent of $\zeta_R(2)$ by Ap\'ery's
theorem.  Within Paper~2's $K = \pi(B + F - \Delta)$, the odd-zeta tier
is a byproduct rather than a summand: $\zeta_R(3)$ is structurally
present on $S^3$ without participating in $K$.

\paragraph{Spinor-intrinsic content from first-order operators on a
spinor bundle.}
The Breit--Pauli reduction of the Dirac equation introduces a second
new transcendental tier: $\alpha^2$ and $\gamma := \sqrt{1 - (Z\alpha)^2}$.
The leading spin-orbit coupling
\begin{equation}\label{eq:bp_so}
  H_{\mathrm{SO}}
  \;=\; \tfrac{Z\alpha^2}{2}\,\tfrac{1}{r^3}\,\mathbf{L}\cdot\mathbf{S},
\end{equation}
evaluated in the $(\kappa, m_j)$ basis via Szmytkowski's reduction
$\langle \mathbf{L}\cdot\mathbf{S}\rangle = -(\kappa+1)/2$ and the
hydrogenic $\langle 1/r^3\rangle_{n,l}
= Z^3/[n^3\, l(l+\tfrac{1}{2})(l+1)]$, gives
\begin{equation}\label{eq:bp_so_matrix}
  \langle n,\kappa,m_j|\,H_{\mathrm{SO}}\,|n,\kappa,m_j\rangle
  \;=\; -\,\frac{Z^4\,\alpha^2\,(\kappa+1)}{4\,n^3\, l(l+\tfrac{1}{2})(l+1)},
\end{equation}
with Kramers cancellation at $l=0$.  Full-Dirac radial corrections enter
through $\gamma$, which satisfies
\begin{equation}\label{eq:gamma_def}
  \gamma^2 + (Z\alpha)^2 \;=\; 1,
\end{equation}
making $\gamma$ algebraic of degree~2 over $\mathbb{Q}(\alpha^2)$.
Together $\alpha^2$ and $\gamma$ generate the ring
\begin{equation}\label{eq:ring_sp}
  R_{\mathrm{sp}} \;:=\; \mathbb{Q}\bigl(\alpha^2\bigr)
    \bigl[\gamma\bigr]/\bigl(\gamma^2 + (Z\alpha)^2 - 1\bigr),
\end{equation}
in which every spinor-intrinsic coefficient lies---mechanically
certified by \texttt{geovac/spinor\_certificate.py}.

This tier is categorically distinct from both the odd-zeta tier and
calibration: $\alpha^2$ is a \emph{coupling} content appearing in
Hamiltonian coefficients (not in the basis or spectrum), depending only
on the existence of spin coupling in the operator.  Higher-order
corrections (Darwin $\alpha^4$, mass-velocity $\alpha^4$, Breit
$\alpha^2$ in $V_{ee}$) populate the same ring $R_{\mathrm{sp}}$ without
creating new transcendental tiers.

\paragraph{Cross-reference to Paper~34's projection taxonomy.}
In the projection-dictionary language of Paper~34 \S~III, the
spinor-intrinsic ring $R_\mathrm{sp}$ is the transcendental signature
of the \emph{Camporesi--Higuchi spinor-lift projection} (Paper~34
\S~III.7), and that projection is the unique source of half-integer
Hurwitz / odd-$\zeta$ / Dirichlet-$\beta$ content in the
fifteen-projection dictionary.  Sprint~TX-A (May~2026, Paper~34
\S~VII Theorem~3 / composition table) confirmed this assignment as
exclusive: no other projection in the catalogue introduces odd-zeta
content into the transcendental signature, and the spinor-lift
projection is the lone bridge between the calibration-tier
($\pi^{\text{even}}$) world of second-order Riemannian operators and
the odd-zeta world of first-order Dirac operators.  Equivalently:
seeing $\zeta_R(3)$, Catalan's $G$, or any Dirichlet-$\beta$ value in
a GeoVac observable is diagnostic of the spinor-lift projection
having been invoked at some step in the chain.


\paragraph{Operator-order $\times$ bundle classification.}

\begin{table}[h]
\caption{Transcendental content classified by operator order and bundle
  type.  The (2nd-order, spinor) cell is degenerate with scalar
  calibration ($\pi^{\mathrm{even}}$ only)---see the $D^2$ discussion
  below.}
\label{tab:bundle_operator_tiers}
\begin{ruledtabular}
\begin{tabular}{llll}
  Tier & Example transcendentals & Operator order & Bundle \\
  \hline
  Intrinsic         & $\pi^{d/2}$ Weyl prefactors & --- & scalar \\
  Calibration       & $\kappa = -1/16$, $F = \pi^2/6$
                    & 2nd ($\Delta_{\text{LB}}$)
                    & scalar \\
  Embedding         & $e^a E_1(a)$, $1/r_{12}$
                    & ---
                    & graph $\to$ $\mathbb{R}^3$ \\
  Flow / composition & $\mu(\rho,R)$, PK
                    & ---
                    & multi-geometry \\
  Odd-zeta          & $\zeta_R(3), \zeta_R(5), \dots$
                    & 1st ($D$)
                    & Dirac on curved space \\
  Spinor-intrinsic  & $\alpha^2$, $\gamma$
                    & 1st ($D$ / Breit--Pauli)
                    & spinor bundle \\
  $D^2$ calibration & $\pi^{\text{even}}$ (degenerate)
                    & 2nd ($D^2$)
                    & spinor bundle \\
\end{tabular}
\end{ruledtabular}
\end{table}

\noindent Tier~1 ($\zeta$) enters through the \emph{spectrum} of the
Dirac operator; Tier~2 ($\alpha^2$, $\gamma$) enters through the
\emph{coupling} of spin to orbital angular momentum in the Hamiltonian
coefficients.  Neither introduces $\pi$ beyond the calibration tier.

\paragraph{Operator-order as transcendental discriminant.}
The squared Dirac operator $D^2$ on unit $S^3$ fills the last open
cell of Table~\ref{tab:bundle_operator_tiers}.  We promote the
discriminating observation to a theorem.

\begin{theorem}[Operator-order discriminant on $S^3$]
\label{thm:operator_order}
Let $D$ be the Camporesi--Higuchi Dirac operator on unit~$S^3$
with spectrum $|\lambda_n| = n + 3/2$ and degeneracy
$g_n^{\mathrm{Dirac}} = 2(n+1)(n+2)$.  Then:
\begin{enumerate}
\item The spectral zeta of the squared operator $D^2$ admits the
exact closed form
\begin{equation}\label{eq:zeta_d2}
  \zeta_{D^2}(s)
  \;=\; 2^{2s-1}\bigl[\lambda(2s{-}2) - \lambda(2s)\bigr],
  \qquad s \in \mathbb{Z}_{\ge 2},
\end{equation}
where $\lambda(a) = (1 - 2^{-a})\,\zeta_R(a)$ is the Dirichlet lambda
function.  Every value lies in $\mathbb{Q}[\pi^2]$ ($\pi^{\mathrm{even}}$
content exclusively).
\item The first-order spectral zeta $D(s) = \sum_n
g_n^{\mathrm{Dirac}}\,|\lambda_n|^{-s}$ admits the Hurwitz
decomposition
\begin{equation}\label{eq:hurwitz_discriminant}
  D(s) \;=\; 2\bigl(2^{s-2}-1\bigr)\,\zeta_R(s{-}2)
           - \tfrac{2^s - 1}{2}\,\zeta_R(s),
  \qquad s \in \mathbb{Z}_{\ge 2}.
\end{equation}
At even~$s$ both terms are $\pi^{\mathrm{even}}$; at odd $s \ge 3$
the terms are $\mathbb{Q}$-linearly independent of $\pi$
(Ap\'ery~1978 for $\zeta_R(3)$; Goncharov--Zagier framework for
higher odd zetas).  In particular $D(4) = 6\zeta_R(2) - \tfrac{15}{2}\zeta_R(4)
= \pi^2 - \pi^4/12$ is $\pi^{\mathrm{even}}$, while the first odd-zeta
witness is $D(5) = 14\,\zeta_R(3) - \tfrac{31}{2}\zeta_R(5)$, carrying
odd-zeta content not present in any value of $\zeta_{D^2}$.
\end{enumerate}
Operator order---first vs.\ second---is therefore the primary
transcendental discriminant on round~$S^3$.  Bundle type modulates
degeneracies, eigenvalue lattices, and rational coefficients, but
not the transcendental class.
\end{theorem}

\begin{proof}[Proof sketch]
For~(1): Lichnerowicz~\cite{lichnerowicz1963} gives $D^2 = \nabla^*
\nabla + R/4$ with scalar curvature $R/4 = 3/2$, so $D^2$ has a
rational shift relative to the conformal Laplacian.  The
Camporesi--Higuchi spectrum yields $|\lambda_n|^2 = (n+3/2)^2$,
which after the half-shift identity for the Hurwitz zeta reduces to
Eq.~(\ref{eq:zeta_d2}).  Bernoulli's formula $\zeta_R(2k) =
\mathrm{rational} \times \pi^{2k}$ then closes the case.  For~(2):
regrouping the same Camporesi--Higuchi sum without squaring keeps the
half-integer shift, producing the Hurwitz combination of
Eq.~(\ref{eq:hurwitz_discriminant}); squaring maps odd-power spectral
sums to even-power sums, structurally eliminating the $\zeta_R(2k{+}1)$
visible in the unsquared sum and reducing the (2nd-order,~spinor) cell
to scalar calibration.  Computational verification:\ \texttt{tests/test\_qed\_vacuum\_polarization.py}
checks Eq.~(\ref{eq:zeta_d2}) symbolically and against a direct
Camporesi--Higuchi spectral sum, and \texttt{tests/test\_qed\_two\_loop.py}
verifies the odd-zeta content of Eq.~(\ref{eq:hurwitz_discriminant}) via
the Hurwitz closed form and a PSLQ parity classification.
\end{proof}

\paragraph{Three-tier collapse (corollary).}
As a corollary of Theorem~\ref{thm:operator_order}, the
operator-order $\times$ bundle taxonomy of
Table~\ref{tab:bundle_operator_tiers} collapses to three effective
tiers:
\begin{enumerate}
  \item \textbf{First-order operators:} odd-zeta content
    ($\zeta_R(3), \zeta_R(5), \ldots$) plus spinor coupling content
    ($\alpha^2$, $\gamma$).
  \item \textbf{Second-order operators:} $\pi^{\text{even}}$ exclusively,
    on both scalar and spinor bundles.
  \item \textbf{Basis and composition tiers:} embedding, flow,
    composition constants (operator-independent).
\end{enumerate}
This is consistent with Paper~24's Coulomb/HO
asymmetry~\cite{loutey_paper24}: second-order Riemannian operators
introduce $\pi$ through nonlinear conformal projections, while
first-order operators (the HO Bargmann Euler operator on
$H^2(S^5)$, the Dirac operator on $S^3$) have $\pi$-free spectra
in rational arithmetic.

\paragraph{Motivic weight and operator order.}
The operator-order discriminant---second-order operators produce
$\pi^{\text{even}}$ content, first-order operators produce odd-zeta
content---has a natural interpretation in the theory of mixed Tate
motives over $\mathrm{Spec}(\mathbb{Z})$.  In this framework
(Goncharov~\cite{goncharov1999}, Zagier~\cite{zagier1994}), the Riemann
zeta value $\zeta_R(n)$ is a \emph{period} of the mixed Tate motive
$\mathfrak{M}(n) = \mathbb{Q}(-n)$ of weight~$n$.  Even zeta values
$\zeta_R(2k) = \text{rational} \times \pi^{2k}$ are periods of motives
that factor through the de~Rham cohomology of even-dimensional compact
manifolds (the Bernoulli formula).  Odd zeta values $\zeta_R(2k{+}1)$
are periods of genuinely independent motives: they do not reduce to
powers of~$\pi$, and $\zeta_R(3)$ in particular is irrational (Ap\'ery
1978) and conjectured algebraically independent from~$\pi$.

The GeoVac operator-order discriminant maps onto this motivic
partition:
\begin{center}
\begin{tabular}{llll}
  \hline
  Operator order & Motivic weight & Transcendental class & Example \\
  \hline
  2nd ($\Delta_{\text{LB}},\; D^2$)
    & even ($2k$) & $\zeta_R(2k) = \text{rat} \cdot \pi^{2k}$
    & $\kappa$, $\mathrm{Vol}(S^3)$, $a_2$ \\
  1st ($D$)
    & odd ($2k{+}1$) & $\zeta_R(2k{+}1) \perp \mathbb{Q}(\pi)$
    & D3: $\zeta_R(3)$ in Dirac Dirichlet \\
  \hline
\end{tabular}
\end{center}

At motivic weight~3, $\zeta_R(3)$ appears as a period of
$\mathbb{Q}(-3)$ in multiple independent geometric contexts:
the Dirac Dirichlet series on~$S^3$ (this work,
Eq.~\ref{eq:dirac_zeta3_witness}), hyperbolic 3-manifold volumes
(Goncharov~\cite{goncharov1999}), two-loop QED $g{-}2$
(Petermann~\cite{petermann1957}; Sommerfield~\cite{sommerfield1957}),
and Chern--Simons theory / algebraic K-theory
($K_5(\mathbb{Z})$ regulator).  The common thread is~$S^3$ or its
quotients.  Zagier's dimension conjecture~\cite{zagier1994} gives
$d_3 = 1$: $\zeta_R(3)$ is the unique independent transcendental
at weight~3, consistent with its single appearance in the GeoVac
Dirac sector.

The connection to QED loop order is structurally predicted by the
motivic correspondence.  The T9 theorem
(Eq.~\ref{eq:zeta_d2}) guarantees that $\zeta_R(3)$ cannot enter through
any single trace of a squared propagator $D^{-2s}$---i.e., at one loop.
It enters at two loops, where irreducible products of first-order
propagators $D^{-1} \otimes D^{-1}$ produce odd motivic weight.  The
operator-order discriminant thus determines the \emph{loop order} at
which new transcendental content first appears: one loop accesses only
even-weight motives ($\pi^{2k}$); two loops access odd-weight motives
($\zeta_R(3)$, $\zeta_R(5)$, \ldots).

The even/odd Hurwitz spectral zeta of the first-order Dirac operator
(Eq.~\ref{eq:hurwitz_discriminant} of Theorem~\ref{thm:operator_order})
provides a sharp computational check.
This exact formula separates the even and odd content at each~$s$:
at $s = 4$ both terms are $\pi^{\text{even}}$
(Eq.~\ref{eq:dirac_d4_value}: $D(4) = \pi^2 - \pi^4/12$); the first
odd-zeta witness on the CH spinor bundle is
Eq.~\eqref{eq:dirac_zeta3_witness}: $D(5) = 14\,\zeta_R(3) -
\tfrac{31}{2}\,\zeta_R(5)$.  At even~$s$ the spectral sum is purely
$\pi^{\text{even}}$; at odd $s \ge 5$, the $\zeta_R(s{-}2)$ and
$\zeta_R(s)$ terms are independently transcendental.  The QED vacuum
polarization at one loop involves the Seeley--DeWitt coefficients
$a_0$, $a_1$, $a_2$ of~$D^2$, which by
Eq.~(\ref{eq:zeta_d2}) are purely $\pi^{\text{even}}$; the
Schwinger~\cite{schwinger1948} result $a_e^{(1)} = \alpha/(2\pi)$
is a single power of~$\pi$, consistent with one-loop = even motivic
weight.  At two loops, $\zeta_R(3)$ appears in the electron $g{-}2$
(Petermann~\cite{petermann1957}; Sommerfield~\cite{sommerfield1957}), consistent with
the discriminant's prediction that odd-weight content requires
irreducible products of first-order propagators.

This motivic interpretation refines the compactness thesis of
Sec.~\ref{sec:intro}: exchange constants are not only projections from
compact to non-compact geometry---they are \emph{periods} of specific
mixed Tate motives, and the operator order on the compact space
determines the motivic weight of the resulting transcendental content.

\paragraph{Dirichlet $L$-values from vertex topology.}
The operator-order discriminant classifies free-propagator content.
A third axis emerges when interaction vertices are included.  The
QED vertex on~$S^3$ couples Dirac spinor harmonics at levels $n_1$,
$n_2$ to a transverse vector harmonic at level~$n_\gamma$ with
SO(4) selection rules requiring $n_1 + n_2 + n_\gamma$ odd (the
$\gamma^\mu$ coupling flips parity).  When this selection rule is
imposed on the two-loop sunset diagram, it splits the Dirac Dirichlet
series into even- and odd-$n$ sub-sums with structurally different
transcendental content.  At $s = 4$ (the packing exponent), the
exact Hurwitz decomposition gives:
\begin{align}
  D_{\text{even}}(4) &= \tfrac{\pi^2}{2} - \tfrac{\pi^4}{24}
    - 4G + 4\beta(4), \label{eq:d_even4} \\
  D_{\text{odd}}(4) &= \tfrac{\pi^2}{2} - \tfrac{\pi^4}{24}
    + 4G - 4\beta(4), \label{eq:d_odd4}
\end{align}
where $G = \sum_{n \ge 0} (-1)^n/(2n{+}1)^2 \approx 0.9159$ is
Catalan's constant and $\beta(4) = \sum_{n \ge 0} (-1)^n/(2n{+}1)^4
\approx 0.9889$ is the Dirichlet beta function at~4.  Both are values
of the $L$-function $L(s, \chi_{-4})$ for the unique non-principal
character mod~4.  Crucially, $G$ and $\beta(4)$ cancel in the full sum
$D_{\text{even}} + D_{\text{odd}} = D(4) = \pi^2 - \pi^4/12$---the
Dirichlet content is invisible without vertex restriction.  The
mechanism is the Hurwitz quarter-shift: even-$n$ modes involve
$\zeta(s, 3/4)$ and odd-$n$ modes $\zeta(s, 5/4)$, whose difference
is $4^s \beta(s)$ (a Dirichlet $L$-value), while their sum reduces
to Riemann zeta ($\pi^{\text{even}}$).

We promote this to a theorem.

\begin{theorem}[Three-axis classification of $S^3$ spectral
transcendentals]
\label{thm:three_axis}
The transcendental content of any spectral functional formed
from the Camporesi--Higuchi Dirac operator $D$, the conformal
scalar Laplacian $\Delta_{\mathrm{LB}}$, and a vertex coupling
on round~$S^3$ is determined by three independent axes:
\begin{enumerate}
\item[(A1)] \emph{Operator order} (first vs.\ second).
By Theorem~\ref{thm:operator_order}, second-order operators
($\Delta_{\mathrm{LB}}$, $D^2$) produce $\pi^{\mathrm{even}}$
content exclusively; first-order operators ($D$) produce odd-zeta
content $\zeta_R(2k{+}1)$ that is $\mathbb{Q}$-linearly independent
of $\pi$.
\item[(A2)] \emph{Bundle type} (scalar vs.\ spinor).
Modulates rational coefficients, eigenvalue lattices, and
degeneracies within the same transcendental class fixed by
axis~A1, but does not change the class.  The (2nd-order, spinor)
cell is degenerate with scalar calibration.
\item[(A3)] \emph{Diagram topology} (free vs.\ vertex-restricted).
Vertex parity selection rules from $\gamma^\mu$ couplings
($n_1 + n_2 + n_\gamma$ odd) project the Dirichlet sum onto
quarter-integer Hurwitz shifts $\zeta(s, 3/4)$, $\zeta(s, 5/4)$
whose difference is $4^s \beta(s) = 4^s L(s, \chi_{-4})$.  This
introduces Dirichlet $L$-values---Catalan's constant
$G = \beta(2)$, $\beta(4)$, and higher---that are absent from
unrestricted free-propagator sums.
\end{enumerate}
Axes A1, A2, A3 are independent in the sense that no axis can be
replaced or absorbed by the others: A1 is computational (operator
$\to$ class), A2 is representation-theoretic (bundle $\to$
coefficient), A3 is interaction-theoretic (vertex
$\to$ character).  Output: the cell at coordinates
$(\mathrm{order}, \mathrm{bundle}, \mathrm{topology})$ determines
the transcendental class.
\end{theorem}

\begin{proof}[Proof sketch]
A1 is Theorem~\ref{thm:operator_order}.  A2 follows from the
$D^2$ closed form (Eq.~\ref{eq:zeta_d2}), where the spinor degeneracy
$g_n^{\mathrm{Dirac}} = 2(n+1)(n+2)$ produces only rational
coefficient shifts relative to the scalar $g_n^{\mathrm{scalar}} =
(n+1)^2$ counterpart---both yield $\pi^{\mathrm{even}}$ at every
integer~$s \ge 2$.  A3 follows from the Hurwitz quarter-shift
identities $\zeta(s, 5/4) - \zeta(s, 3/4) = -4^s \beta(s)$ applied
to the vertex-parity-restricted sum, giving
Eqs.~(\ref{eq:d_even4})--(\ref{eq:d_odd4}); the unrestricted sum is
$D_{\mathrm{even}}(s) + D_{\mathrm{odd}}(s) = D(s)$, in which $G$
and $\beta(4)$ cancel exactly, so the Dirichlet content is invisible
without vertex restriction.  Computational verification: PSLQ at
80-digit precision; 35/35 tests in
\texttt{tests/test\_qed\_vertex.py}.
\end{proof}

\paragraph{Remark: Dirichlet $L$-tier as new transcendental class.}
The Dirichlet $\beta$-values introduced by axis~A3 are a genuinely
new transcendental class: they arise from the \emph{interaction}
(how electron and photon modes couple at the vertex), not from the
free propagators.  In the exchange-constant taxonomy of
Sec.~\ref{sec:taxonomy} they constitute a \emph{vertex exchange
constant} tier living beyond both calibration~$\pi$ and Riemann
odd-zeta.  This is the entry that the RH-sprint $\chi_{-4}$ identity
(Sec.~\ref{sec:level5_rh_sprint}, Eq.~\ref{eq:rh_j_identity}) places
in closed form for the full Dirac Dirichlet difference.

\paragraph{Photon vector structure as a calibration exchange constant.}
The graph-native QED construction (Paper~28, GN-1 through GN-7)
computes one-loop QED entirely on the finite Fock graph using exact
algebraic arithmetic.  On the graph, the photon lives on edges
(1-cochains) as a \emph{scalar} object---the edge Laplacian
$L_1 = B^T\!B$ propagates it without reference to any vector
structure.  In the continuum, the physical photon is a transverse
vector field on~$S^3$ carrying $\mathrm{SO}(4) = \mathrm{SU}(2)_L
\times \mathrm{SU}(2)_R$ quantum numbers, and this vector structure
enforces selection rules (vertex parity, SO(4) channel counting~$W$)
that are absent on the graph.  Adding explicit photon angular momentum
labels $(q, m_q)$ with Clebsch--Gordan vertex coupling and E1 parity
enforcement recovers 7~of~8 continuum QED selection rules for
\emph{scalar} electrons (Paper~28).  The sole failure for scalar
electrons is Furry's theorem: diagonal vertex couplings
$V(a,a,q,m_q) \neq 0$ for $l \geq 1$ when $q$ is odd.  For
\emph{Dirac} electrons, Furry's theorem is recovered as well,
achieving 8/8 (see below).

The calibration cost of the $1 \to 7$ transition is precisely one
transcendental per loop: $\pi$, entering as the $S^2$ Weyl exchange
constant $1/(4\pi)$ from the vector spherical harmonic normalization
$\sqrt{(2l+1)(2q+1)(2l'+1)/(4\pi)}$.  This is a concrete instance
of the calibration exchange constants cataloged in the taxonomy table
(Table~\ref{tab:taxonomy}): the $S^2$ Weyl constant $\omega_2/(2\pi)^2
= 1/(4\pi)$ converts the discrete angular momentum quantum numbers
into the continuous solid angle measure required for vector photon
normalization.  The entire scalar graph QED is $\pi$-free
($\mathbb{Q}[\sqrt{2}, \sqrt{3}, \sqrt{6}, \ldots]$); the vector
photon introduces exactly the $\omega_2 = \pi$ entry from the
Weyl constant table.

In the taxonomy, the photon's angular momentum structure therefore
sits in the \textbf{calibration} tier, alongside the volume
normalization $\mathrm{Vol}(S^3) = 2\pi^2$ and the Hurwitz-zeta
regularization of infinite spectral sums: it is a property of the
\emph{manifold embedding}, not of the graph.  This sharpens the
graph/continuum partition for QED: the graph captures the
algebraic/combinatorial content (rational propagators, algebraic CG
vertices, $\pi$-free observables), while the continuum projection
adds three calibration-tier ingredients: (1)~volume normalization,
(2)~infinite-sum regularization, and (3)~photon vector structure
(cost: $1/(4\pi)$ per loop).
The 8~continuum selection rules partition by transcendental cost
into three tiers.  For \emph{scalar} electrons, the partition is
$1$--$7$--$8$: graph topology (1~rule, zero cost)
$\to$ angular momentum conservation (6~rules, calibration cost
$1/(4\pi)$) $\to$ field theory (1~rule, requires second-quantized
$C$-symmetry).  For \emph{Dirac} electrons, the partition refines
to $1$--$6$--$1$--$8$: graph topology (1~rule, intrinsic, zero cost)
$\to$ angular momentum conservation (6~rules, calibration $1/(4\pi)$)
$\to$ Dirac spinor phase constraint (1~rule, intrinsic, zero cost).

\paragraph{The Dirac spinor phase constraint as an intrinsic tier.}
When the electron is described by a Dirac spinor
$\psi = (f \cdot \Omega_\kappa,\; i g \cdot \Omega_{-\kappa})^T$
(Paper~23), the $\alpha$-matrix vertex coupling is off-diagonal in the
large/small component space.  The factor of~$i$ on the small component
makes the diagonal coupling $V(a,a,q,m_q)$ vanish identically for all
states: $V_{LS} + V_{SL} = 0$ by Hermiticity.  This is the mechanism
by which Furry's theorem is recovered---tadpole diagrams (self-loops at
vertex~$a$) vanish because the single-particle Dirac vertex is
kinematically off-diagonal.

The key taxonomic point is the \emph{cost} of this recovery.  The
$1 \to 7$ transition (adding photon angular momentum labels) costs
exactly $\pi$ per loop---it is calibration content.  The $7 \to 8$
transition (Dirac spinor phase) costs \emph{nothing}: no new
transcendental enters, no continuum embedding is required.  The spinor
phase constraint is a kinematic identity of the single-particle Dirac
vertex, not a field-theoretic symmetry ($C$-conjugation).  It lives in
the same \textbf{intrinsic} tier as the graph topology rule (Gaunt/CG
sparsity), not in the calibration tier.

This distinction has a concrete implication: the $7 \to 8$ transition
does \emph{not} require second-quantized field theory.  Any
single-particle framework using the Dirac equation on~$S^3$ with
vector-photon coupling automatically recovers all~8 selection rules.
The standard attribution of Furry's theorem to charge conjugation
symmetry is operationally correct in flat-space QFT but obscures the
fact that, on the compact $S^3$ graph, the same selection rule has a
simpler origin in the Dirac spinor structure itself.

\paragraph{Flat-space consistency of the transcendental classification.}
The three-axis taxonomy is derived from $S^3$ spectral sums at finite
radius.  A natural concern is whether the classification is a
finite-curvature artifact.  We verify that it is not.
On $S^3$ of radius~$R$, the Dirac eigenvalues scale as
$|\lambda_n(R)| = (n + 3/2)/R$ while the degeneracies $g_n$ are
$R$-independent.  The Dirichlet series therefore satisfies
$D(s,R) = R^s\,D(s,1)$ \emph{exactly}, making the ratio
$D(s,R)/R^s$ an $R$-independent structural invariant.  In the
$R \to \infty$ limit, Weyl's law confirms that the mode count
$N(n_{\max}) = 2(n_{\max}+1)(n_{\max}+2)(n_{\max}+3)/3$
matches the flat-space momentum-integral density
$(2/3)\,R^3 p_{\max}^3$, verifying the spectral-to-momentum
correspondence.  The even/odd $s$-parity discriminant persists at
every~$R$: $D(s_{\text{even}},R)/R^s$ is $\pi^{\text{even}}$;
$D(s_{\text{odd}},R)/R^s$ contains odd-zeta.  The \emph{same}
transcendental constants---$\zeta_R(3)$, $G$, $\beta(4)$---that
appear in the $S^3$ spectral sums also appear in flat-space
two-loop QED (Petermann~\cite{petermann1957};
Sommerfield~\cite{sommerfield1957}), confirming that curvature
modifies rational coefficients but not the transcendental class.
An honest limitation: matching the full Petermann--Sommerfield
coefficient $a_2 = -197/144 + \pi^2/12 + 3\zeta_R(3)/4
- \pi^2 \ln 2 / 2$ requires diagram-specific computation
(renormalization subtractions, SO(4) Clebsch--Gordan algebra)
beyond the scope of structural spectral sums.
The flat-space limit thus validates the taxonomy as a
\emph{structural} property of the operator algebra, independent
of the $S^3$ compactification scale.

\paragraph{Irreducibility of the Level~4 piecewise structure.}
The piecewise-smooth $\rho$-dependence of $\mu(\rho,R)$ at Level~4
is not a coordinate artifact: it cannot be resolved by geometric
elevation to any higher-dimensional space (Track~Z).  Three approaches
were ruled out: (i)~algebraic blow-up of $|s - \rho|$ produces smooth
sheets but with $\rho$-dependent integration limits ($\alpha = \arccos\rho$,
$\arcsin\rho$), making matrix elements transcendental in~$\rho$;
(ii)~Lie algebra elevation fails because $\min(s,\rho)/\max(s,\rho)$
is $C^0$ but not~$C^1$, incompatible with matrix elements of any
finite-dimensional Lie group representation (which are real-analytic);
(iii)~$S^3 \times S^3$ product space fails because two-center nuclear
attraction cannot be smooth on any single compact manifold.
The piecewise structure thus belongs to the embedding category
in the taxonomy: it is the irreducible geometric content of placing
two nuclei at finite separation in hyperspherical coordinates, analogous
to the $1/r_{12}$ codimension obstruction (Track~W) but affecting the
nuclear rather than the electronic coupling.

\paragraph{Even-odd staircase and channel convergence rate.}
An extended convergence study at Level~4 (Paper~15, Track~AA)
reveals a pronounced even-odd staircase in $D_e$ vs $l_{\max}$:
even increments produce large gains ($+42$, $+6.8$, $+1.7$~pp)
while odd increments produce near-zero gains ($+0.66$, $+0.08$~pp).
The root cause (Track~AC) connects directly to the exchange constant
structure: the nuclear coupling matrix
$\langle l_1' l_2' | V_{\rm nuc} | l_1 l_2 \rangle$ is diagonal
in one electron's angular momentum index (Kronecker delta), so
direct coupling from the ground channel $(0,0)$ requires a
$\Delta l = 2$ transition enforced by the gerade constraint
$l_1 + l_2 = \text{even}$.  Odd-$l_{\max}$ channels couple only
indirectly through the electron-electron multipole expansion.

This selection rule constrains the \emph{convergence rate} of the
partial-wave expansion: the nuclear coupling---a piecewise exchange
constant in the taxonomy---couples to the ground state only through
quadrupolar ($\Delta l = 2$) transitions, producing an algebraic
convergence rate $\sim (l + 1/2)^{-2}$ with an even-odd modulation.
The $\sigma$ and $\pi$ channel sectors are completely decoupled
(both nuclear and $e$--$e$ coupling conserve $m$), with $\pi$
contributing a constant ${\sim}6.6$~pp additive offset.

\paragraph{PK pseudopotential as composition exchange constant.}
The Phillips--Kleinman pseudopotential in the composed geometry
(Paper~17) enforces core-valence orthogonality, which is a
Pauli exclusion effect---distinct from the Coulomb exchange
constants cataloged above.  Without PK, no equilibrium
geometry exists (the PES is monotonically attractive), establishing
PK as essential for creating the bonding minimum.  Ab initio
derivation of the $R$-dependent PK scaling parameter $R_{\rm ref}$
from core properties fails due to a fundamental scale mismatch:
all core-derived candidates ($0.27$--$0.82$~bohr) are too small
by a factor of $4$--$11$.

A complete survey of all solver $\times$ PK combinations
(Paper~17, Table~VI) confirms that PK is not an improvable
approximation.  The 2D variational solver eliminates ${\sim}75\%$
of the $l_{\max}$ drift (from the adiabatic approximation), but
$l$-dependent PK has zero additional effect on the 2D solver.
The residual drift ($+0.15$~bohr/$l_{\max}$) is intrinsic to
the composed-geometry factorization.

Three graph-native alternatives to PK were investigated and all
failed: (i)~node exclusion between the core $S^3$ and valence
Level~4 graphs has no coordinate correspondence;
(ii)~the fiber bundle Berry connection is intra-group, and the
exchange operator $P_{ij}$ mixes coordinate systems;
(iii)~spectral orthogonality (function-space) does not imply
coordinate-space exclusion.  The master obstruction is that
antisymmetry requires a shared coordinate system; composition
factorizes $\Psi_{\rm total}$ into incompatible coordinates.

PK therefore defines a sixth exchange type in the taxonomy:
\textbf{composition constants} convert between a single-geometry
(full $N$-electron) and a composed-geometry (factorized) description.
Unlike the other exchange constants, which convert between discrete
and continuous descriptions, composition constants measure the
accuracy cost of the factorization itself.  The ``exchange rate''
is quantified by the 144$\times$ angular compression that composed
LiH achieves over a full 4-electron solver on $S^{11}$
(Paper~17, Sec.~V): the ${\sim}5\%$ $R_{\rm eq}$ error at
$l_{\max} = 2$ is the irreducible cost of this compression.

This classification is validated by the full 4-electron Level~4N
solver (Paper~17, Sec.~VIII): exact $S_4$ antisymmetry produces
molecular equilibrium without PK at $l_{\max} \geq 2$, confirming
that PK approximates---rather than creates---the physics of Pauli
exclusion.  The composition exchange constant is the cost of
replacing exact inter-group antisymmetry with a pseudopotential
in the factorized coordinate system.

\paragraph{Inner-factor input data: a fourth tier disjoint from
the outer-factor taxonomy.}
The exchange constants cataloged above all live on the GeoVac
\emph{outer} factor: the Fock-projected $S^3$ graph, its spinor
lift, and their continuum limits.  When the outer triple is
combined with a Connes--Chamseddine \emph{inner} finite spectral
triple $T_F = (A_F, H_F, D_F, J_F, \gamma_F)$ to form an
almost-commutative product $T = T_{\mathrm{GV}} \otimes T_F$
(Paper~32 \S~VIII.C; Sprint~H1), the master Mellin engine
$\mathcal{M}[\mathrm{Tr}(D^k\,e^{-tD^2})]$ extends multiplicatively
to the combined triple but with two structural reductions on the
inner side.

\begin{theorem}[$\eta$-trivialization on Krajewski-class finite
triples]
\label{thm:eta_trivialization}
Let $(A_F, H_F, D_F, J_F, \gamma_F)$ be a finite even spectral triple
in the Krajewski--Paschke--Sitarz lineage~\cite{krajewski1998,
paschke_sitarz2000}.  The chirality grading axiom
$\{\gamma_F, D_F\} = 0$ implies, for every odd $k \in \mathbb{N}$,
\begin{equation}\label{eq:eta_trivial}
  \mathrm{Tr}\bigl(D_F^k \cdot e^{-t D_F^2}\bigr) \;\equiv\; 0,
  \qquad t > 0.
\end{equation}
In particular, the M3 sub-mechanism (vertex-parity Hurwitz, $k=1$)
of the master Mellin engine vanishes identically on any
Connes--Chamseddine-compatible inner factor.
\end{theorem}

\begin{proof}
$\{\gamma_F, D_F\} = 0$ and $\gamma_F^2 = 1$ make $\gamma_F$ a unitary
involution conjugating $D_F$ to $-D_F$.  Conjugation preserves
traces, so $\mathrm{Tr}(D_F^k\,e^{-tD_F^2}) = \mathrm{Tr}((-D_F)^k
e^{-tD_F^2}) = (-1)^k \mathrm{Tr}(D_F^k\,e^{-tD_F^2})$.  Odd $k$
forces zero. $\square$
\end{proof}

\begin{theorem}[Multiplicative factorization of the combined
spectral action]
\label{thm:ac_factorization}
For the combined Dirac $D = D_{\mathrm{GV}} \otimes 1_F + \gamma_{\mathrm{GV}}
\otimes D_F$ with $\gamma_{\mathrm{GV}}$ the outer chirality
grading,
\begin{equation}\label{eq:d_squared_factorize}
  D^2 \;=\; D_{\mathrm{GV}}^2 \otimes 1_F + 1_{\mathrm{GV}}
            \otimes D_F^2,
\end{equation}
the cross term $\{D_{\mathrm{GV}} \otimes 1, \gamma_{\mathrm{GV}}
\otimes D_F\} = (D_{\mathrm{GV}}\gamma_{\mathrm{GV}} +
\gamma_{\mathrm{GV}}D_{\mathrm{GV}}) \otimes D_F$ vanishing by
outer chirality anticommutation.  Consequently
\begin{equation}\label{eq:trace_factorize}
  \mathrm{Tr}\bigl(e^{-tD^2}\bigr)
  \;=\; \mathrm{Tr}\bigl(e^{-tD_{\mathrm{GV}}^2}\bigr)
        \cdot \mathrm{Tr}\bigl(e^{-tD_F^2}\bigr),
\end{equation}
and the master Mellin engine on the combined triple factorizes as
(outer M$_i$ ring)~$\times$~(inner Yukawa Dirichlet ring).
\end{theorem}

The Mellin output of the inner factor is a generalized Dirichlet
series in the eigenvalues $\{y_i\}$ of $D_F$ (the Yukawa parameters,
in the SM realization $A_F = \mathbb{C} \oplus \mathbb{H} \oplus
M_3(\mathbb{C})$):
\begin{equation}\label{eq:inner_dirichlet}
  \mathcal{M}\bigl[\mathrm{Tr}(D_F^k \cdot e^{-t D_F^2})\bigr](s)
  \;=\; \Gamma(s) \cdot \sum_i m_i \cdot y_i^{\,k - 2s},
  \qquad k \in \{0, 2\},
\end{equation}
with rational multiplicities $m_i$ (including the SU(3) color
factor~3 for the quark sectors).  This output lies in the ring
$\mathbb{Q}\bigl[y_1^{-2s}, y_2^{-2s}, \ldots, y_n^{-2s}\bigr]$
---a \emph{parameter-tied Dirichlet ring} whose generators are
free inputs to the framework, not transcendentals derived from any
outer-factor mechanism.

We propose a sixth named tier in the taxonomy:

\begin{quote}
\textbf{Inner-factor input data}: parameter-tied Dirichlet rings
generated by the eigenvalues of a finite Dirac operator $D_F$ on
a Connes--Chamseddine-compatible inner factor.  Source: the choice
of finite spectral triple $T_F$ and its Yukawa data.  Type:
empirical input from outside the framework, categorically disjoint
from the intrinsic, conformal/calibration, embedding,
algebraic-implicit, and composition tiers of the outer factor.
\end{quote}

This tier sharpens the ``no autonomous SM-distinguishing data''
verdict of Sprints H1, G3, and G4a (Paper~32 \S~VIII.C): the
Yukawa parameters, hypercharges, and mixing angles of the SM
inhabit a ring with no overlap with any outer-factor ring, and
GeoVac's taxonomy provides no derivation principle for them.
The framework controls the outer factor's transcendental content
via M1/M2/M3; the inner factor controls the SM-distinguishing
parameters via Eq.~(\ref{eq:inner_dirichlet}); the two tensor
cleanly by Theorem~\ref{thm:ac_factorization}, and neither
determines the other.

The structural separation is consistent with the
universal/Coulomb partition of Paper~31~\cite{loutey_paper31} read
at the spectral-triple level: the universal sector (algebra $A$,
selection rules) transfers freely; the potential-specific sector
(Dirac operator $D_{\mathrm{GV}}$, calibration $\kappa$, $\alpha$
content) requires per-system rederivation; and now the inner
sector (algebra $A_F$, finite Dirac $D_F$) requires empirical
specification of its eigenvalue spectrum.  Three structural
sources of input data, three categorically disjoint tiers.

Computational evidence: the four-leg investigation (driver
\texttt{debug/finite\_spectral\_triple\_engine.py}, memo
\texttt{debug/inner\_factor\_mellin\_engine\_memo.md}) verified
Theorems~\ref{thm:eta_trivialization}--\ref{thm:ac_factorization}
symbolically on the SM lepton sector, the full SM
$A_F = \mathbb{C} \oplus \mathbb{H} \oplus M_3(\mathbb{C})$, and
two comparators ($\mathbb{C} \oplus \mathbb{C}$ and
$M_2(\mathbb{C}) \oplus M_3(\mathbb{C})$).  The KO-dim~6 cut
on Krajewski candidates (required to complete the outer KO-dim~3
to the SM-compatible total~9~$\equiv$~1~$\pmod 8$) and
$\eta$-trivialization together leave a large candidate
sub-class; the engine does \emph{not} reverse-engineer $A_F$.
This is a clean negative on the ``second packing axiom''
question: GeoVac's machinery says where $A_F$'s Dirichlet ring
sits in the taxonomy but does not pick out which finite spectral
triple realizes it.

\paragraph{Open question:\ a second packing axiom for inner-factor data.}
The Paper~0 packing axiom forces the outer spectral triple's
discrete labels $(n, l, m, s)$ and the $S^3$ topology, but is silent
on inner-factor content $A_F$.  Sprint H1's $\mathcal{A}_{\mathrm{GV}}
\otimes (\mathbb{C} \oplus \mathbb{H})$ extension produces a Higgs
sector for any non-zero Yukawa~$Y$, but no GeoVac data couples to
that selection (the Sprint G3 commuting-$\mathbb{Z}_2$ theorem
$\| \gamma_{\mathrm{GV}} \otimes I_F - I_{\mathrm{GV}} \otimes
\gamma_F \|_{\mathrm{op}} = 2$ is the structural fingerprint).
The Bertrand $\times$ Hopf-tower truncation provides only an upper
bound on gauge content;\ the lower-bound forcing (why $\mathrm{SU}(2)$
and $U(1)$ appear alongside $\mathrm{SU}(3)$, why generation count $= 3$,
why specific Yukawa ratios) is calibration-side.  Whether a
\emph{sibling} axiom---distinct from the Paper~0 packing axiom---generates
this data is an open question of the framework.  Two candidate routes
(charge-universal generation map;\ $\mathrm{SU}(5) / E_8 / \mathrm{Spin}(10)$
topological embedding) were tested in Phase~4H and falsified
(Tracks SM-B, SM-C; CLAUDE.md~\S 3).  The honest scope statement is
that GeoVac's machinery places inner-factor data in this taxonomic
tier but does not derive it.

\paragraph{Thermal probes of the Yukawa Dirichlet ring.}
Sprint TD Track~3 (May 2026,
\verb|debug/sprint_td_track3_memo.md|) catalogues five thermal
observables whose values depend on parameters of the inner-factor
Yukawa Dirichlet ring of Eq.~(\ref{eq:inner_dirichlet}): lepton
thermal free energy ($\Sigma\,y_i^2$ at $k=2$); QCD thermal pressure
$g_{\mathrm{QCD}}$ ($M_3(\mathbb{C})$ multiplicity ratio at $k=0$);
Higgs effective potential at finite $T$ (off-diagonal CC block
trace at $k=2$); electroweak Debye masses ($(2 + N_{\mathrm{gen}})\,g^2$
at $k=2$); and thermal running $\alpha(T)$ (rational charge sum and
Yukawa-tied logarithmic cutoffs).  Each observable's outer baseline
factors cleanly through the Track~1 tensor-product spectral triple
$T_{S^3}\otimes T_{S^1_\beta}$ (Sprint TD Track~1; Paper~35
\S~\ref{sec:tensor_product_verification}), with the Stefan--Boltzmann
$M_1\times M_2 = -\pi^2/90\,T^4$ identity as the universal backbone.
Each inner residual reduces to a generalized Dirichlet sum
$\sum_i c_i\, y_i^{-2s}$ ($k=0$) or $\sum_i c_i\, y_i^{2-2s}$ ($k=2$)
at integer~$s$ in the Yukawa eigenvalues.  Three further candidates
(axial chemical potential thermodynamics, parity-odd transport,
finite-$T$ EDMs) are filtered out by Theorem~\ref{thm:eta_trivialization}:
they reduce to $k=1$ traces on the inner factor and identically
vanish, a structural prediction of the framework.  None of the five
surviving observables are predicted by GeoVac alone---they are
consistency checks that constrain candidate $T_F$, not derivations.
The headline empirical content is that the Marcolli--van~Suijlekom
gauge-network reading without Higgs ($Y = 0$) predicts no high-$T$
electroweak symmetry restoration and is empirically falsified by the
sphaleron transition at $T_C \approx 160$~GeV; this sharpens the
positive-thin H1 verdict.

\paragraph{Chemistry-side analog: external atomic-physics calibration
data.}
The same structural classification applies to GeoVac's composed
natural-geometry chemistry side.  The composed-qubit architecture
builds heavy-atom molecules ($Z \ge 11$) using a frozen-core
treatment: the [Ne], [Ar], [Kr], [Xe] cores are reduced to a screened
$Z_{\rm eff}(r)$ profile via the FrozenCore solver, with hydrogenic
plus multi-zeta valence orbitals built on top.  The five
Clementi--Raimondi 1963 exponents per row, the multi-zeta
coefficients, and the screened-core orbital shapes are external
atomic-physics inputs to the framework---they are not derived from
any outer-factor Mellin mechanism (M1/M2/M3).

Sprint W1e period-class diagnostic (2026-06-04;
\verb|debug/sprint_w1e_period_class_memo.md|) tested whether the
W1c\,$\to$\,W1e chemistry-correlation wall sits in a higher-weight
cyclotomic-mixed-Tate (M3) period class than the diagonal of the
composed Hamiltonian.  Eleven NaH correction terms extracted from
sprints F4 / F5 / F6 were PSLQ-classified at 100~dps against M1, M2,
M3, and an INNER analog basis built from the Clementi--Raimondi
exponents.  Result: zero outer-factor (M1/M2/M3) identifications across
all eleven terms at audit ceiling 100 and permissive ceiling $10^6$,
after filtering for rational-only fits, basis-internal $\mathbb{Q}$-linear
dependencies, and the curve-fit-audit failure mode (basis size
vs.\ input precision).  An independent random-rational null-hypothesis
test against the same bases gave $0/50$ false positives on M2 and M3,
confirming the bases are sparse enough that ``0~hits across 11~targets''
is informative, not a measurement-noise artifact.

The PSLQ result is a consistency check, not a derivation: the
load-bearing argument is structural.  The M3 sub-mechanism lives at
$k=1$ in the master Mellin engine
$\mathcal{M}[\mathrm{Tr}(D^k\,e^{-tD^2})]$ (Sec.~\ref{sec:mellin}),
which ties to vertex-parity / parity-character spectral sums on the
GeoVac graph.  The W1e correction terms have no vertex-parity
content---they are FCI eigenvalues and Hartree integrals over
hydrogenic plus multi-zeta orbital products---so they categorically
cannot pick up M3 content.  The PSLQ pass rules out the loophole
``W1e might happen to land in M1/M2/M3 by coincidence''; the
structural argument is what carries the negative.

The chemistry-side corrections classify cleanly into the chemistry-side
analog of the inner-factor input-data tier: FCI eigenvalues are
algebraic-implicit roots of polynomials with externally-fitted matrix
elements; Hartree $J$ and Slater--Condon $K$ integrals are rational
in external $Z_{\rm eff}$ values; Phillips--Kleinman barriers depend
on the same external orbital fits.  Six W1c / W1e closure attempts
(CLAUDE.md~\S{}3 entries for cross-center PK, screened-Schrödinger
valence, multi-zeta substitution, kernel-shape substitution,
bonding-PK rank-1, basis enlargement, explicit-core Hartree) are
explained at the structural level by this non-selection: they attempt
outer-factor mechanisms (rational/algebraic modifications) on a wall
whose content lies entirely in the external input-data tier.

The chemistry analog of Sprint H1's Yukawa non-selection theorem is
the W1e wall: the framework can compose any Clementi--Raimondi
exponent set into the heavy-atom Hamiltonian but cannot autonomously
select the exponents.  Outer-factor M-mechanism machinery is silent
on inner-factor / calibration-data content in both the SM-side
(Yukawa values) and the chemistry-side (atomic-physics calibration
set) directions.  The five chemistry-side observables in the
Sprint~TD Track~3 analog role (alkali-hydride binding energies,
ionisation energies, fine-structure splittings, hyperfine splittings,
polarisabilities) are not predicted by GeoVac alone; they are
consistency checks that constrain candidate atomic-physics calibration
sets, not framework-derived predictions.

\paragraph{Operational localization (Sprint R3-B falsifier,
2026-06-07).}  The W1e classification at the period-class level
(0/11 hits in M1/M2/M3) is now sharpened by an operational
test.  Sprint R3-B exported the NaH balanced Hamiltonian
to FCIDUMP at $n_{\max} \in \{2, 3\}$ via the
\texttt{GeoVacHamiltonian.to\_fcidump} interface and ran
DMRG (sector-restricted FCI at $\mathrm{FCI\,dim} \le 784$,
which IS DMRG-at-infinite-$\chi$) at every $R$ in
$[2.5, 6.0]$ bohr.  The resulting PES is \emph{bit-identical to
the P4 direct-FCI baseline} ($\max\text{diff} \sim 6 \times
10^{-13}$~Ha at every grid point), with no interior minimum
and $21$--$24\times$ over-binding at the experimental
$R_e = 3.566$ bohr.  Increasing $n_{\max}$ from 2 to 3
\emph{deepens} the unphysical well rather than closing it.
The wall is therefore not at the determinant-expansion level
(DMRG cannot move it) and not at basis-truncation level
($n_{\max}$ sweep deepens it); it lives at the projection
step from continuous Level 4 adiabatic+PK to second-quantized
$(h_1, \mathrm{eri}, e_{\rm core})$ integrals.  A parallel
diagnostic on LiH composed and balanced (Sprint R3-A,
same date) found the same monotone-descending qubit-FCI PES
behaviour: the LiH 2.82\% $R_{\rm eq}$ established at
CHANGELOG v3.56.0 came from
\texttt{ComposedDiatomicSolver.LiH\_ab\_initio} (continuous
Level 4 multichannel adiabatic with PK), \emph{not} from FCI
on the qubit/FCIDUMP-exported integrals.  The
chemistry-side inner-factor input-data tier therefore has a
specific structural home in the GeoVac construction: the
projection step where continuous Level 4 multichannel
information is reduced to the second-quantized integrals that
seed qubit/composed solvers.  This is the chemistry-side
\emph{operational} analog of the Sprint H1 Yukawa
non-selection theorem at the projection level: the framework
admits a chemistry Hamiltonian structurally but does not
autonomously generate the calibration content of its
$(h_1, \mathrm{eri}, e_{\rm core})$ tensors at this
projection step.


%% ========================================================================
%% V. THE FINE STRUCTURE CONSTANT CONNECTION
%% ========================================================================

\section{Connection to the Fine Structure Constant}
\label{sec:alpha}

Paper~2~\cite{loutey_paper2} observes that the fine structure
constant $\alpha$ satisfies the cubic $\alpha^3 - K\alpha + 1 = 0$,
where the coefficient
\begin{equation}
  K = \pi\!\left(42 + \frac{\pi^2}{6} - \frac{1}{40}\right)
  = 137.036064
  \label{eq:K}
\end{equation}
is constructed from spectral invariants of the Hopf fibration
$S^1 \to S^3 \to S^2$.  The combination rule
$K = \pi(B + F - \Delta)$ is recorded as an Observation---an
empirical formula with structural support but no first-principles
derivation.

We observe that this combination rule has precisely the structure
of a spectral-geometric exchange:
\begin{itemize}
  \item $B = 42$ is discrete representation-theoretic data:
    the degeneracy-weighted $\mathrm{SO}(3)$ Casimir trace over
    the first three shells of $S^3$.
  \item $F = \zeta(2) = \pi^2/6$ is a continuous spectral
    invariant: the spectral zeta function of $S^1$ at $s = 1$.
  \item $\Delta = 1/40$ is a boundary correction mixing discrete
    eigenvalue data ($|\lambda_3| = 8$) with the cumulative state
    count ($N(2) = 5$).
  \item The overall factor $\pi$ is $\omega_2 = \pi$, the 2D
    ball volume, which serves as the numerator of the $S^2$ Weyl
    density and as the packing-plane prefactor
    $\sigma_0 = \pi d_0^2 / 2$ in Paper~0 (the base of the Hopf
    fibration, the photon propagation sphere).
\end{itemize}

The structure of $K$ is thus: the 2D ball volume $\omega_2 = \pi$
(the numerator of the $S^2$ Weyl density, and the same $\pi$ that
fixes Paper~0's packing-plane prefactor $\sigma_0 = \pi d_0^2 / 2$)
multiplying a sum of the $S^2$ Casimir trace ($B$, discrete),
the $S^1$ spectral zeta value ($F$, transcendental), and an
$S^3$ boundary term ($\Delta$, mixed).  This is a Weyl-type
formula that converts between the discrete bundle
(representation labels on $S^3$) and the continuous coupling
(the electromagnetic interaction strength measured in
$\mathbb{R}^3$).

If this identification is correct, then the combination rule
$K = \pi(B + F - \Delta)$ is not arbitrary numerology but an
instance of the same spectral-geometric exchange that produces
$\kappa = -1/16$ at Level~1 and $e^a E_1(a)$ at Level~2.
The exchange constant for the Hopf bundle is $K$ itself---the
``cost'' of projecting the $S^3$ electron state space through
the $S^1$ gauge fiber onto the $S^2$ photon base, measured
in units of coupling strength.

\begin{observation}[Spectral-geometric reading]
The combination rule $K = \pi(B + F - \Delta)$ has the shape of a
spectral-geometric exchange constant for the Hopf fibration,
analogous to the Weyl--Selberg constants for compact manifolds.
Whether the specific form---$\pi$ multiplying a sum of Casimir data
and zeta values---admits a derivation from the spectral geometry
of $S^1 \to S^3 \to S^2$ (e.g.\ via the heat-kernel expansion or the
analytic torsion of the Hopf bundle) is an open question; no such
derivation is known. The combination rule itself remains an Observation.
\end{observation}

Two additional observations support this identification:

\begin{observation}[Structural match]
The combination rule $K = \pi(B + F - \Delta)$ mixes discrete
data ($B = 42$, from representation theory) with continuous
spectral data ($F = \zeta(2) = \pi^2/6$, from the $S^1$
spectral zeta function), multiplied by an overall geometric
factor ($\omega_2 = \pi$, the 2D ball volume that enters as the
numerator of the $S^2$ Weyl density and as the packing-plane
prefactor $\sigma_0 = \pi d_0^2 / 2$ in Paper~0, with $S^2$ the
Hopf base).  This is precisely the pattern identified throughout this
paper: a transcendental conversion factor mediating between
discrete and continuous descriptions of the same geometry.
The structural match is evidence---not proof---that the
combination rule is an instance of the exchange constant pattern
rather than numerology.
\end{observation}

\begin{observation}[$\alpha$ is likely transcendental]
The fine structure constant $\alpha$ is a root of the cubic
$\alpha^3 - K\alpha + 1 = 0$, where $K$ involves $\pi$ and
$\zeta(2) = \pi^2/6$.  Roots of polynomials with transcendental
coefficients \textit{can} be algebraic (e.g., $x^2 - \pi x +
\pi(\pi - 1) = 0$ has the algebraic root $x = 1$), but this
requires special cancellations that are generically absent.
Since $K$ is a non-trivial combination of $\pi$ and $\pi^2$
with rational coefficients, it is overwhelmingly likely
(though not proven) that $\alpha$ is transcendental---consistent
with its role as a spectral-geometric exchange constant.
\end{observation}

\subsection{Three-tier structure of the combination rule}
\label{sec:alpha_three_tier}

An ongoing computational investigation into the origin of the three
components of $K = \pi(B + F - \Delta)$---using exact-rational
symbolic computation and high-precision arithmetic---has pinned down
$B$ and $F$ as genuinely different objects on the Fock $(n,l)$
lattice, and has eliminated a broad family of candidate mechanisms
for $\Delta$. We observe that the three terms appear to occupy
three categorically distinct tiers of the exchange-constant
taxonomy rather than specializing a single common generator. We
summarize the four structural findings below; none of them alters
the Observation-level status of Paper~2's cubic identity.

\paragraph{1. A $\kappa$--$B$ identity from the Fock weight.}
The Fock projection of the hydrogenic basis onto $S^3$ carries a
residual momentum-space weight
$w(p) = (p^2 + p_0^2)^{-2}$
after the overall $(2p_0)^{5/2}$ prefactor of the Fock map is
stripped.  Computation reveals that the diagonal matrix elements
$\langle n,l|w|n,l\rangle_{p_0 = 1}$ for $n \le 3$ are exact
rationals whose denominators divide $32 = 2\,d_{\max}^2$:
\begin{equation}
  \{\langle n,l|w|n,l\rangle_{p_0 = 1}\}_{(n,l)}
  = \Bigl\{ \tfrac{7}{16},\, \tfrac{5}{8},\, \tfrac{3}{8},\,
            \tfrac{5}{8},\, \tfrac{17}{32},\, \tfrac{11}{32} \Bigr\}
  \label{eq:fock_weight_values}
\end{equation}
for $(n,l) \in \{(1,0),(2,0),(2,1),(3,0),(3,1),(3,2)\}$.  Weighted
by the Hopf base factor $(2l+1)\,l(l+1)$ and summed,
\begin{equation}
  \sum_{n=1}^{3}\sum_{l=0}^{n-1} (2l+1)\,l(l+1)\,
    \langle n,l|w|n,l\rangle_{p_0 = 1}
  = \frac{63}{4}
  = 6\,B\,|\kappa|,
  \label{eq:kappa_B_identity}
\end{equation}
with $B = 42$ the Casimir trace of Paper~2 and $\kappa = -1/16$
the calibration constant of Sec.~\ref{sec:kappa_derivation}.
Equivalently, averaging over the six $(n,l)$ cells gives
$B\,|\kappa| = 21/8$.  This is the first exact rational algebraic
link in the GeoVac framework between the calibration constant and
the Hopf base invariant: $|\kappa| = 1/d_{\max}^2$ appears as the
universal denominator factor of the Fock-weight matrix elements,
while $B$ appears in the Hopf-weighted sum.  The identity
(\ref{eq:kappa_B_identity}) is exact-rational with no fitted parameters
(the arithmetic closes:\ $9/4 + 51/16 + 165/16 = 63/4 = 6\,B\,|\kappa|$);
the six Fock-weight values of Eq.~(\ref{eq:fock_weight_values}) are
stated here, and a regression test deriving them from the graph and
checking the identity is a logged coverage gap
(\texttt{docs/claim\_test\_matrix.md}).

\paragraph{2. A Dirichlet identity for $F$.}
Define the Dirichlet series of the $S^3$ Fock shell degeneracy
$g_n = n^2$,
\begin{equation}
  D_{n^2}(s) \;:=\; \sum_{n=1}^{\infty} g_n\, n^{-s}
             \;=\; \sum_{n=1}^{\infty} n^{-(s-2)}
             \;=\; \zeta_R(s - 2),
  \label{eq:Dn2_def}
\end{equation}
evaluated at the packing exponent $s = d_{\max} = 4$ of Paper~0:
\begin{equation}
  D_{n^2}(d_{\max}) = \zeta_R(2) = \frac{\pi^2}{6} = F.
  \label{eq:F_dirichlet}
\end{equation}
The exponent $s = 4$ admits three independent natural readings:
$s = d_{\max}$ from Paper~0's packing axiom; $s = \dim(\mathbb{R}^4)$,
the ambient Euclidean space into which $S^3$ is stereographically
projected in Paper~7; and $s = 2 N_{\mathrm{init}}$ from the same
packing axiom.  The weight $n^2$ is the degeneracy of the $n$-th
$S^3$ shell in Fock's projection.  To our knowledge
(\ref{eq:F_dirichlet}) is the first identification of $F$ from a
graph-intrinsic quantity in the GeoVac framework---that is, $F$ is
produced here without invoking an embedding metric or a
sphere-spectral construction.  No other natural Dirichlet series
on the shell lattice (per-shell Casimir, state count, Laplacian
eigenvalue, $l$-max weight) produces an exact hit at any integer
$s$ in the range $[3, 12]$.

\paragraph{3. Three independent spectral homes (Theorem 3).}
Combining the $\kappa$--$B$ identity (\ref{eq:kappa_B_identity})
with the Dirichlet identity (\ref{eq:F_dirichlet}) and the
Paper~2 canonical form
$1/\Delta = |\lambda_{n_{\max}}|\, N(n_{\max}{-}1)$, we promote the
three-tier structure to a theorem; the additive combination rule
itself remains an Observation.

\begin{theorem}[Three independent spectral homes for $B$, $F$, $\Delta$]
\label{thm:k_decomposition}
Let $K = \pi(B + F - \Delta)$ be the combination of
Paper~2~\cite{loutey_paper2} with $B = 42$, $F = \pi^2/6$,
$\Delta = 1/40$.  The three components have structurally independent
spectral origins on the Fock-projected $S^3$ graph:
\begin{enumerate}
\item[(B)] $B = \sum_{n=1}^{3} \sum_{l=0}^{n-1} (2l{+}1)\,l(l{+}1)
= 42$ is a \emph{finite} Casimir trace of the scalar
$\mathrm{SO}(3)$ Laplacian, truncated at $n_{\max} = 3$;
algebraic class: rational integer.
\item[(F)] $F = D_{n^2}(s = d_{\max}) = \zeta_R(2) = \pi^2/6$
is the \emph{infinite} Dirichlet series of the Fock-shell degeneracy
$g_n = n^2$ at the packing exponent $s = d_{\max} = 4$ of Paper~0;
algebraic class: transcendental, in $\mathbb{Q}\cdot\pi^2$.
\item[($\Delta$)] $\Delta^{-1} = |\lambda_{n_{\max}}|\,
N(n_{\max}{-}1) = 8 \cdot 5 = g_{n=3}^{\mathrm{Dirac}} = 40$ is the
\emph{boundary} product at the truncation edge---equivalently the
single-chirality Camporesi--Higuchi Dirac mode degeneracy at
$n_{\mathrm{CH}} = 3$; algebraic class: rational, finite, on the
\emph{spinor} bundle.
\end{enumerate}
The three objects do not arise from a single common generator:
$B$ is a finite combinatorial sum, $F$ is an infinite arithmetic
series, $\Delta$ is a finite boundary degeneracy on a different
bundle (spinor vs.\ scalar).  Each lives in a distinct cell of the
operator-order $\times$ bundle classification of
Theorem~\ref{thm:operator_order} ($B$: 2nd-order / scalar /
finite-truncation; $F$: 2nd-order / scalar / infinite-Dirichlet;
$\Delta$: 1st-order / spinor / boundary).
\end{theorem}

\begin{proof}[Proof sketch]
For~$(B)$ and~$(F)$: Eqs.~(\ref{eq:kappa_B_identity}) and
(\ref{eq:F_dirichlet}) above, by direct exact-rational computation.
For~$(\Delta)$: the Camporesi--Higuchi spectrum
$|\lambda_n| = n + 3/2$ with degeneracy
$g_n^{\mathrm{Dirac}} = 2(n+1)(n+2)$ gives $g_3^{\mathrm{Dirac}} = 40$
at the Paper~2 selection-principle cutoff $n_{\mathrm{CH}} = 3$.
The scalar-Laplacian factorization $|\lambda_3|\cdot N(2) =
8 \cdot 5 = 40$ is a different presentation of the same integer; the
spinor Dirac form $g_3^{\mathrm{Dirac}}$ is a single polynomial
identity $2(n{+}1)(n{+}2)|_{n=3} = 40$, structurally cleaner.
Independence: paragraph~4 below catalogues the negative results
that exclude any common Dirichlet, sphere-spectral, or Hopf-fiber
generator simultaneously producing $B$, $F$, and $\Delta$.
\end{proof}

\begin{observation}[Combination rule, restated]
\label{obs:k_rule}
The additive combination $K = \pi(B + F - \Delta)$, with the three
components characterized by Theorem~\ref{thm:k_decomposition},
equals $\alpha^{-1}$ to $8.8 \times 10^{-8}$.  No first-principles
derivation of the combination is known.  In the lattice gauge
language of Paper~25~\cite{loutey_paper25} and the non-abelian SU(2)
extension of Paper~30~\cite{loutey_paper30}, $K/\pi$ has the shape
of a finite-cutoff spectral-action coefficient (matter trace +
gauge-fiber zeta + Dirac boundary mode count), but the equality with
$\alpha^{-1}$ is an empirical observation, not a theorem.
The supertrace analysis on unit~$S^3$ (Paper~28~\cite{loutey_paper28},
Sprint~2 Q-tracks) confirms that all three components are invisible
to the perturbative Connes--Chamseddine asymptotic expansion---$K/\pi$
lives in the non-perturbative sector if it lives in any
spectral-action sector at all.
\end{observation}

\paragraph{4. Negative results on candidate mechanisms.}
The three-tier reading is supported by the systematic elimination
of alternative unifying mechanisms.  Computation (exact rational
symbolic arithmetic and \texttt{mpmath} at 50 decimal places)
has ruled out, within the tolerance indicated:
(i) the Hopf-twist spectral comparison $S^3$ versus $S^1 \times S^2$
(cleanest near-miss $\sim 6.8 \times 10^{-3}$);
(ii) higher Casimir traces on $S^3$ at all orders through
$n_{\max} = 3$;
(iii) Hopf graph quotient Laplacian invariants for
$S^3 \to S^2$ over the $m$-fiber;
(iv) the discrete $S^1$ fiber spectral zeta (path-graph Laplacian,
in four weighting schemes);
(v) the continuous $S^1$ fiber via Mellin and heat-kernel expansions
(the Jacobi inversion $\theta_{S^1}(t) \sim \sqrt{\pi/t}$ forces every
order to be linear in $\pi$: the $\sqrt{\pi}\cdot\sqrt{\pi}\cdot\mathbb{Q}
= \pi\cdot\mathbb{Q}$ structure collapses any expected $\zeta_R(2)$
contribution to $\pi$ times a rational at every order);
(vi) $S^5$ spectral geometry, including zetas, Seeley--DeWitt
coefficients, and functional determinants (all rational, or
reducible to $\zeta_R(\mathrm{odd})$ and $\log 2$);
and (vii) a $\zeta$-combination scan for $\Delta$ over $3{,}228$
candidates built from Riemann and Hurwitz zetas, Bernoulli numbers,
and Stieltjes constants, with zero strict hits within $10^{-25}$
of $1/40$.

The structural reason for these negative results is that
round-sphere Laplace--Beltrami operators produce (i) integer powers
of $\pi$ in their volumes, (ii) rational Seeley--DeWitt
coefficients, and (iii) $\zeta_R(\mathrm{odd}) + \log 2$ combinations
in their spectral determinants---but never $\pi^2 = 2\zeta_R(2)$
additively next to a rational.  This is consistent with the
$\pi$-free certificate of the Bargmann--Segal lattice
(Paper~24), which shows that the
analogous lattice for the 3D harmonic oscillator on the Hardy
space of $S^5$ is bit-exactly rational in exact arithmetic; and the
same rigidity holds for $S^5$ Laplace--Beltrami spectral invariants
generally.

\paragraph{Reframed open question.}
The derivation question ``where does each piece of $K$ come from?''
has been partially answered: $B$ from a finite Casimir truncation
(with the $\kappa$--$B$ link of Eq.~\ref{eq:kappa_B_identity}), $F$
from an infinite Dirichlet series at the packing exponent
(Eq.~\ref{eq:F_dirichlet}), and $\Delta$ from a boundary product at
the truncation edge (irreducible to $\zeta$-content within the
candidate set tested).  The remaining structural question is no
longer a derivation question about any single component: it is a
question about why the \emph{additive} combination
$B + F - \Delta$, multiplied by the overall factor $\omega_2 = \pi$,
equals $\alpha^{-1}$ to $8.8 \times 10^{-8}$.  That is a question
about the conspiracy of three categorically different objects---a
finite sum, an infinite Dirichlet series, and a boundary
product---rather than the specialization of a single generator.
Paper~2's identification remains an Observation.


%% ========================================================================
%% VI. OBSERVABLE CLASSIFICATION
%% ========================================================================

\section{Observable Classification by Transcendental Content}
\label{sec:observable_classification}

The exchange constant taxonomy (Sec.~\ref{sec:taxonomy}) classifies the
\emph{constants themselves} by their mathematical origin.  A complementary
classification addresses the \emph{observables}: given a physical quantity
to be computed, what class of transcendental content does it require?

The conformal equivalence between the graph and continuous descriptions
(Paper~7) implies a natural partition of observables into three classes,
determined by what type of projection is needed to evaluate them.

\begin{enumerate}
  \item \textbf{Class~S (spectral observables):} Energy eigenvalues,
    degeneracies, selection rules, and transition amplitudes between
    eigenstates.  These require only the graph eigenvalue spectrum
    and the calibration constant $\kappa$.  No transcendental content
    beyond $\kappa = -1/16$ (rational) is needed.

  \item \textbf{Class~P (single-particle spatial observables):}
    Electron density $\rho(\mathbf{r})$, expectation values
    $\langle r^k \rangle$, momentum distributions, and any quantity
    requiring the wavefunction evaluated at a point in $\mathbb{R}^3$.
    These require the conformal projection $S^3 \to \mathbb{R}^3$,
    which introduces $\pi$ through the volume element
    $d\Omega_{S^3} = 2\pi^2$ and the spherical harmonic
    normalization $\int |Y_{lm}|^2 d\Omega = 1$ (with
    $d\Omega$ carrying $4\pi$).

  \item \textbf{Class~C (multi-particle correlated observables):}
    Correlation energies, adiabatic potential energy surfaces,
    multi-electron expectation values, and any quantity requiring
    the inter-particle coordinate.  These require the higher exchange
    constants: $e^a E_1(a)$ at Level~2 for $m \neq 0$ states
    (Sec.~\ref{sec:stieltjes}), $\mu(R)$ at Level~3
    (Sec.~\ref{sec:mu}), $\mu(\rho,R)$ at Level~4
    (Sec.~\ref{sec:mu_level4}), and the PK composition constant
    at Level~5.
\end{enumerate}

\noindent This classification is stated as a structural claim:

\medskip
\noindent\textbf{Claim~4.}
\textit{The class of transcendental content required to evaluate a
quantum-mechanical observable is determined by the type of projection
linking the graph description to the coordinate system in which the
observable is defined.  Spectral observables (Class~S) require no
transcendental content.  Single-particle spatial observables (Class~P)
require the Weyl constants ($\pi^{d/2}$).  Multi-particle correlated
observables (Class~C) require the embedding, flow, or composition
exchange constants of Sec.~\ref{sec:taxonomy}.}

\medskip

Claim~4 is falsifiable: for any observable, one can determine
its class, identify the required exchange constants, and verify
that no additional transcendental content enters the computation.
The following worked examples illustrate each class.

\subsection{Example 1: Hydrogen energy levels (Class~S)}

The hydrogen energy spectrum $E_n = -Z^2/(2n^2)$ is a spectral
observable.  The continuum $S^3$ Laplace--Beltrami operator that the
graph Laplacian converges to has eigenvalues
$\lambda_n = -(n^2 - 1)$, which are integers for all $n$
(Paper~7, Eq.~(12); Paper~0, Sec.~II).  The degeneracies
$g_n = n^2$ are integers.  The Fock conformal mapping (Paper~7,
Sec.~IV) converts these integer eigenvalues to the Rydberg
spectrum via the energy-shell constraint $p_0^2 = -2E$
and the stereographic focal length:
\begin{equation}
  E_n = -\frac{Z^2}{2n^2}\;\text{Ha},
  \label{eq:hydrogen_class_s}
\end{equation}
with the calibration constant $\kappa = -1/16$ (rational)
mediating between graph eigenvalues and physical energies
(Sec.~\ref{sec:kappa_derivation}).  Every quantity in this
derivation---eigenvalues, degeneracies, $\kappa$---is an
integer or rational number.  No transcendental content enters.

Selection rules provide a second example.  The dipole selection
rule $\Delta l = \pm 1$ follows from the triangle inequality
on the Gaunt integral
$\int Y_{l_1 m_1} Y_{1q} Y_{l_2 m_2}\,d\Omega$, which
reduces to a product of Wigner $3j$ symbols---rational numbers
computed from factorials of integers~\cite{loutey_paper0}.
No transcendental content enters.

\subsection{Example 2: Electron density (Class~P)}

The electron density $\rho(\mathbf{r}) = |\psi(\mathbf{r})|^2$
requires evaluating the wavefunction in position space.  On the
graph, the state $|n, l, m\rangle$ is a node with an integer
label; in $\mathbb{R}^3$, the corresponding wavefunction is
$\psi_{nlm}(\mathbf{r}) = R_{nl}(r)\,Y_{lm}(\theta,\phi)$.
The conformal projection $S^3 \to \mathbb{R}^3$ (Paper~7,
Sec.~III) introduces the conformal factor
$\Omega = 2p_0/(p^2 + p_0^2)$, which mediates between the
unit-sphere measure $d\Omega_{S^3}$ and the flat-space volume
element $d^3\mathbf{r}$.  The spherical harmonic normalization
$\int |Y_{lm}|^2\,d\Omega = 1$ introduces $\pi$ through the
solid angle $\oint d\Omega = 4\pi$.  The radial normalization
$\int_0^\infty |R_{nl}|^2\,r^2\,dr = 1$ introduces no additional
transcendentals for hydrogen (the integrals are rational multiples
of $a_0^{-3}$ times exponentials, evaluated in closed form).

The density is thus a Class~P observable: it requires the $\pi$
that enters through the projection from the dimensionless $S^3$
topology to the dimensionful $\mathbb{R}^3$ coordinate system.

\subsection{Example 3: Helium ground state energy (Class~C)}

The helium ground state energy $E_0 = -2.9037$~Ha~\cite{drake2006}
is a correlated observable.  The Level~3 hyperspherical solver
(Paper~13) computes it via the adiabatic potential $U(R)$, which
depends on the angular eigenvalue $\mu(R)$.  The function $\mu(R)$
satisfies the characteristic polynomial $P(R,\mu) = 0$
(Sec.~\ref{sec:algebraic_curve}), with coefficients in
$\mathbb{Q}(\pi, \sqrt{2})$ inherited from the angular coupling
matrix~\cite{loutey_track_p1}.  The energy is obtained by
integrating the hyperradial Schr\"odinger equation with
$U(R)$ as the effective potential---a definite integral over
the algebraic curve, producing a transcendental number.

On the $S^3 \times S^3$ product space (Level~3 FCI), the same
energy can in principle be computed without $\mu(R)$: the
Hamiltonian matrix has entries given by Gaunt integrals (rational)
and the Slater $F^0$ direct Coulomb integrals (algebraic on $S^3$,
Paper~7 Sec.~VI).  The FCI result (0.35\% error) demonstrates that
the correlated energy \emph{can} be approximated within the
algebraic framework, but the $\mu(R)$ parameterization is needed
to achieve sub-0.1\% accuracy.  The transition from Class~S to
Class~C reflects the physical content: electron-electron
correlation couples the two $S^3$ factors, and the coupling
strength varies with the hyperradial coordinate $R$---a spatial
degree of freedom that exists only in the projected description.

\paragraph{Three-layer decomposition (v2.6.0).}
The algebraic Casimir CI (Track~DI Sprint~2) makes the exchange
constant structure of the He ground state fully explicit.  At each
$n_{\max}$, the FCI matrix $H(k) = Bk + Ck^2$ has rational entries
as a function of the orbital exponent~$k$, so the variationally
optimal energy $E^*$ is an algebraic number---a root of a polynomial
with rational coefficients.  At $n_{\max} = 1$: $k^* = 27/16$,
$E^* = -729/256$~Ha (exact rationals).  This reveals a three-layer
structure:
\begin{center}
\begin{tabular}{lll}
\hline
Layer & Content & Exchange type \\
\hline
Rational & Slater integrals, Gaunt coefficients & None (graph-intrinsic) \\
Algebraic & Optimal exponents at finite $n_{\max}$ & Calibration \\
Transcendental & e-e cusp ($\alpha = \pi/4$ on $S^5$) & Embedding \\
\hline
\end{tabular}
\end{center}
The Fock energy-shell self-consistency $k^2 = -2E$, which is exact
at Level~1 (one electron, one scale), over-constrains the
two-electron problem: self-consistent~$k$ gives 5--13\% error vs.\
1.5--1.9\% for variational~$k$.  This is the first instance where a
single calibration constant is insufficient---the transition from
calibration to embedding exchange constants corresponds to the
physical content of electron-electron correlation.

\subsection{Boundary cases}

Two boundary cases clarify the classification:

\begin{itemize}
  \item \textbf{Ionization energies} are spectral (Class~S) when
    computed as eigenvalue differences $\Delta E = E_n - E_m$,
    which involve only $\kappa$ and integer quantum numbers.
    However, photoionization \emph{cross sections} require the
    dipole matrix element $\langle \psi_f | \mathbf{r} | \psi_i
    \rangle$, which is a spatial observable (Class~P) because
    it evaluates the wavefunction at positions in $\mathbb{R}^3$.

  \item \textbf{Two-electron energies} computed via FCI on
    $S^3 \times S^3$ are formally Class~S (graph eigenvalues
    of the product-space Hamiltonian), but achieve limited
    accuracy due to the product-basis completeness deficiency
    (Sec.~\ref{sec:mu_level4}).  High-accuracy two-electron
    energies require the Class~C exchange constants.
    The classification thus tracks both the type of observable
    and the accuracy target.
\end{itemize}

The classification does not assert that one class is more
fundamental than another.  The conformal equivalence
(Paper~7) ensures that all three descriptions---graph, projected
single-particle, correlated multi-particle---are mathematically
equivalent.  The classification identifies which exchange
constants mediate between them.


%% ========================================================================
%% VII. DISCUSSION
%% ========================================================================

\section{Discussion}
\label{sec:discussion}

\subsection{Transcendence and the graph design principle}

Exchange constants are transcendental because they convert between
incommensurable descriptions: a discrete count (eigenvalues,
representation labels) and a continuous measure (volume, curvature,
geodesic length).  Weyl constants carry $\pi^{d/2}$, and $\pi$ is
transcendental by Lindemann's theorem.  Not all functions in the
hierarchy are transcendental: $\mu(R)$ satisfies a polynomial $P(R,\mu) = 0$
(Sec.~\ref{sec:algebraic_curve}), making it algebraic in~$R$, with
transcendental content inherited from angular coupling
coefficients~\cite{berger2003}.  The irreducible transcendentals are
the \emph{constants} that enter the coupling matrix, not the
parameterized functions built from them.

When the computation stays on the graph, exchange constants are
unnecessary.  This suggests a design principle: \textit{minimize
coordinate projections} to minimize transcendental content.  The
ideal solver stays on the graph; every departure incurs a
transcendental toll.

\subsection{Exchange constants are computationally cheap}

Conceptual complexity does not correlate with computational cost:
$\kappa$ is one multiplication; $e^a E_1(a)$ is one function
evaluation seeding an $O(N^2)$ recurrence; $\mu(R)$ is a
$10$--$30$ dimensional matrix diagonalization per $R$-point.
The hierarchy tracks curvature mismatch severity, not computational
difficulty.  At Level~2, identifying the Stieltjes seed and absorbing
it into the basis \emph{reduces} cost (eliminating quadrature) while
improving convergence (Sec.~\ref{sec:stieltjes}).

\subsection{The hierarchy of transcendence}

Exchange constants increase in complexity with curvature mismatch:
constant ($\kappa = -1/16$, Level~1) $\to$ single transcendental
seed ($e^a E_1(a)$, Level~2) $\to$ algebraic function of one
parameter ($\mu(R)$ satisfying $P(R,\mu) = 0$ over
$\mathbb{Q}(\pi,\sqrt{2})$, Level~3) $\to$ piecewise-algebraic
function of two parameters ($\mu(\rho,R)$, Level~4).  This mirrors
the mathematical hierarchy: Weyl constants are periods
($\pi = \int dx/(1+x^2)$), Stieltjes seeds are exponential periods
($E_1(a) = \int_1^\infty e^{-at}/t\,dt$), and adiabatic eigenvalues
are algebraic functions of the continuous parameter built from
transcendental coefficients.


\subsection{Structural parallels}

The pattern parallels string compactification, where compact internal
geometry determines low-energy couplings~\cite{polchinski1998}.
We note also that
$\mathrm{SO}(4,2) = \mathrm{Iso}(\mathrm{AdS}_5)$, the isometry group
of five-dimensional anti-de~Sitter space, is the conformal group of the
boundary $\mathbb{R}^{3,1}$; the Fock projection
$S^3 \to \mathbb{R}^3$ is a Euclidean section of this conformal
boundary correspondence.

\subsection{The algebraic curve refinement}
\label{sec:algebraic_curve}

The characteristic polynomial analysis~\cite{loutey_track_p1}
refines the exchange constant taxonomy by distinguishing between
transcendental \emph{constants} and algebraic \emph{functions} of
continuous parameters:
\begin{itemize}
  \item \textbf{Truly irreducible transcendentals:}
    The embedding constants $e^a E_1(a)$ (Level~2, $m \neq 0$) and the
    angular coupling integrals $c/\pi$ and $c\sqrt{2}/\pi$
    (Level~3 and~4) are irreducible---they cannot be eliminated by
    any algebraic manipulation.  These are the genuine ``exchange
    constants'' in the Weyl--Selberg sense.
  \item \textbf{Algebraic functions of transcendental data:}
    The adiabatic eigenvalues $\mu(R)$ and $\mu(\rho,R)$ are
    algebraic functions whose coefficients inherit transcendentals
    from the coupling matrix.  The $R$-dependence (or $(\rho,R)$-dependence)
    is algebraic at every truncation: $\mu$ satisfies a polynomial
    equation $P(R,\mu) = 0$ of degree $d$ equal to the matrix dimension.
    At $l_{\max} = 0$, $P$ is linear and $\mu(R)$ is an explicit
    rational function of $R$ (with transcendental coefficients).
    At $l_{\max} = L$, $P$ has degree $L+1$ in both $R$ and $\mu$.
  \item \textbf{Calibration constant:} $\kappa = -1/16$ is rational
    (Sec.~\ref{sec:kappa_derivation}).
\end{itemize}

This refinement has a practical implication: the ``transcendental toll''
of the adiabatic parameterization is \emph{inherited}, not \emph{accrued}.
Moving from the product space $S^3 \times S^3$ to hyperspherical
coordinates does not introduce new transcendental content; it
rearranges existing transcendentals (from the angular coupling)
into an algebraic function of $R$.  The toll is paid once, in the
coupling matrix, and the $R$-dependence is algebraically determined.

This does not yet eliminate the computational cost of pointwise
diagonalization: at $l_{\max} \geq 1$, the characteristic polynomial
has degree $\geq 2$ in $\mu$ and must be solved numerically at each
$R$-point.  But it does mean that the solver operates on a known
polynomial, not an opaque transcendental function.  Whether this
algebraic structure can be exploited computationally---for instance,
via resultant methods, homotopy continuation, or algebraic
PES representations---is an open question.

This algebraic structure is specific to Level~3, where the angular
Hamiltonian $H(R) = H_0 + R \cdot V_C$ is a linear matrix pencil.
At Level~4, the split-region Legendre expansion of the nuclear charge
function introduces piecewise-smooth $\rho$-dependence that breaks
the pencil structure (Track~S).  The Level~4 angular eigenvalues
are piecewise-algebraic but do not satisfy a global characteristic
polynomial.

\subsection{Two routes to the value of $\kappa = -1/16$}
\label{sec:kappa_derivation}

The original identification of $\kappa = -1/16$ (Paper~0) was
empirical: the unique constant that maps GeoVac graph eigenvalues
onto the Rydberg energy levels.  Below we give two independent routes
to its numerical value and classify their agreement.  As recorded in
Paper~0 and Paper~7, that the calibration scale coincides with a
conformal-geometry quantity is an \emph{Observation}---a numerical
coincidence with no known derivation bridging the two---not a
derivation of $\kappa$'s calibration role from the geometry.

\textit{Graph Laplacian spectrum.}  The GeoVac lattice
(Paper~0) discretizes $S^3$ with states $|n,l,m\rangle$
connected by angular transitions $m \leftrightarrow m \pm 1$
and radial transitions $n \leftrightarrow n \pm 1$.  Interior
nodes have degree $d_{\max} = 4$ (two angular, two radial
neighbors).  The graph Laplacian $\mathcal{L} = D - A$ is
positive semidefinite; its maximum eigenvalue satisfies
$\lambda_{\max} \leq 2 d_{\max} = 8$, with the bound saturated
in the continuum limit ($n_{\max} \to \infty$).  Numerical
verification confirms monotonic convergence:
$\lambda_{\max}(n_{\max}{=}5) = 6.62$,
$\lambda_{\max}(10) = 7.61$,
$\lambda_{\max}(20) = 7.90 \to 8$.

\textit{Energy calibration.}  The graph Hamiltonian
$H = \kappa \mathcal{L}$ (no node weights) has its ground
state at energy $\kappa \lambda_{\max}$, since $\kappa < 0$
maps the largest graph eigenvalue to the lowest physical energy.
Matching this to the hydrogen ground state:
\begin{equation}
  \kappa \lambda_{\max} = E_H = -\tfrac{1}{2}
  \implies
  \kappa = \frac{-1/2}{8} = -\frac{1}{16}.
  \label{eq:kappa_derived}
\end{equation}

\textit{Factor identification.}  The value $-1/16 = -(1/2)(1/8)$
decomposes into:
\begin{enumerate}
  \item $1/2$: the non-relativistic kinetic energy normalization
    $E = p^2/(2m)$, or equivalently the energy-shell constraint
    $E = -p_0^2/2$;
  \item $1/8 = 1/(2 d_{\max})$: the inverse of the maximum
    graph Laplacian eigenvalue, determined by the vertex degree
    of the lattice discretization.
\end{enumerate}

\textit{Heat kernel test.}  The Seeley--DeWitt coefficients on unit
$S^3$ ($a_0 = 2\pi^2$, $a_1 = 2\pi^2$, $a_2 = \pi^2$) produce no
ratio matching $1/16$; neither does the Weyl counting function
$N(\Lambda) \sim \Lambda^{3/2}/3$ or the spectral zeta.  The conformal
coupling $c_3 = 1/8$ for $d=3$ numerically equals $1/\lambda_{\max}$,
but this is coincidental: $c_3$ is a Yamabe invariant while
$1/\lambda_{\max}$ depends on graph vertex degree.

\textit{Fock weight derivation (Note added v2.26.1).}  The same value
is derivable independently from the Fock projection itself, without
reference to graph vertex degree.  The $l = 0$ eigenfunctions on $S^3$
are Gegenbauer polynomials $C^1_{n-1}(\cos\chi)$ (equivalently,
Chebyshev polynomials of the second kind $U_{n-1}(\cos\chi)$).  Their
three-term recurrence gives transition amplitudes of $1/4$ between
adjacent shells, so the squared Fock coupling is
\begin{equation}
  c^2(n,l) = \tfrac{1}{16}\bigl[1 - l(l{+}1)/(n(n{+}1))\bigr],
  \label{eq:fock_coupling_p18}
\end{equation}
with $c^2(n,0) = 1/16$ universally for all $n$.  Equivalently,
$1/16 = 1/\Omega^4(0)$ where $\Omega(0) = 2$ is the stereographic
conformal factor at $p = 0$.  The graph-degree derivation
(Eq.~\ref{eq:kappa_derived}) and the Fock weight route yield the same
value:\ the graph's maximum vertex degree $d_{\max} = 4$ and the
conformal factor $\Omega(0) = 2$ are both fixed by the unit-$S^3$
geometry the graph discretizes ($2 d_{\max} = \Omega(0)^4 = 2^4/2 = 8$).
That the calibration scale $\kappa$ equals this geometric value is the
Observation;\ no derivation links the two.

A notable special case: $c^2(4,3) = 1/40 = \Delta$ (Paper~2), linking
the highest-$l$ Fock coupling at the Paper~2 cutoff to the Dirac
degeneracy $g_3^{\mathrm{Dirac}} = 40$.

\textit{Classification.}  In Table~\ref{tab:taxonomy} $\kappa$ sits in
the ``conformal'' tier:\ its numerical value $1/16$ \emph{coincides}
with the conformal-geometry quantity $1/\Omega^4(0)$ of the Fock
projection, a value any discretization of $S^3$ would share.  The two
routes (graph degree and conformal factor) return the same number
because both are computed from the unit-$S^3$ geometry;\ that the
calibration scale $\kappa$ equals it is an \emph{Observation} (a
numerical coincidence, no known derivation bridge), consistent with the
Observation-tier status of $\kappa$ recorded in Paper~0 and Paper~7.  The Z-universality noted in
Paper~0 ($\kappa \to \kappa Z^2$ for any nuclear charge)
reflects the Z-independence of the graph structure, with only
the energy-shell scaling $E \propto Z^2$.


%% ========================================================================
%% VIII. STRUCTURAL INCOMMENSURABILITY (META-PATTERN)
%% ========================================================================

\section{Structural Incommensurability within GeoVac Observables}
\label{sec:meta_observation}

The exchange-constant taxonomy developed in the preceding sections
has repeatedly encountered a distinctive pattern: single framework
observables in GeoVac, when unpacked, resolve into sums or
products of categorically distinct transcendental or spectral
objects with \emph{no common generator}.  The pattern has been
observed in three independently-derived three-axis classifications.
We name it for future reference: \emph{structural incommensurability
within GeoVac observables}.

\begin{observation}[Structural incommensurability]
\label{obs:incommensurability}
Across the framework, single-observable statements
(``two-loop QED spectral sum on $S^3$'',
``$K = \pi(B + F - \Delta)$'',
``zero structure of the GeoVac zeta'')
have repeatedly resolved under structural analysis into sums or
superpositions of objects from categorically different cells of the
operator-order $\times$ bundle $\times$ topology classification of
Theorem~\ref{thm:three_axis}.  The objects at each cell are what they
are, not derivations of each other; the appropriate taxonomy is
three-axis, not single-generator.  The status is empirical at three
documented instances; it may sharpen if further framework observables
decompose the same way or if a bridge is found that fails cleanly.
\end{observation}

\paragraph{Instance~1 (QED on $S^3$, Paper~28).}  Perturbative QED
sums on~$S^3$~\cite{loutey_paper28} produce three transcendental
classes simultaneously: even-zeta content ($\pi^{\mathrm{even}}$,
from second-order Riemannian operators), odd-zeta content
($\zeta(3), \zeta(5)$, from first-order Dirac), and Dirichlet-$L$
content ($G = \beta(2), \beta(4)$, from vertex topology).  These
three classes appear together in the two-loop sunset sum but are
discriminated by Theorem~\ref{thm:three_axis}'s three axes
(operator order $\times$ bundle type $\times$ vertex parity), and
no single mechanism generates all three.

\paragraph{Instance~2 ($\alpha$-combination decomposition).}  By
Theorem~\ref{thm:k_decomposition}, the observed combination rule
$K = \pi(B + F - \Delta)$ has three components with structurally
distinct spectral homes: $B = 42$ is a finite Casimir trace on the
scalar Laplacian at $m = 3$; $F = \pi^2/6$ is the infinite Dirichlet
series of Fock degeneracies at the packing exponent
$d_{\max} = 4$; and $\Delta^{-1} = 40$ is the degeneracy
$g_3^{\mathrm{Dirac}} = 2 \cdot 4 \cdot 5$ of the third Dirac mode on
the spinor bundle.  No common generator produces all three; $K$
assembles three structurally-independent spectral objects.  In the
gauge-structure language of Paper~25~\cite{loutey_paper25} and the
SU(2) Wilson extension of Paper~30~\cite{loutey_paper30}, $K/\pi$
has the shape of a finite-cutoff spectral-action coefficient
(matter trace + gauge zeta + Dirac boundary mode count).

\paragraph{Instance~3 (Weil dictionary on the GeoVac Hopf graphs,
Paper~29).}  The Ihara-zeta side (combinatorial, closed
non-backtracking walks) and the spectral-Dirichlet side
(Dirac-eigenvalue Dirichlet series) of the same
Hopf graph~\cite{loutey_paper29} give \emph{opposite} verdicts.
Zero spacing: Ihara Poisson ($\mathrm{CV} \approx 1$), spectral
GUE ($\mathrm{CV} \approx 0.35$).  Dirichlet-$L$ content: Ihara
absent (bipartiteness makes the $\chi_{-4}$ support disjoint from
walk-length support); spectral present in the closed form of
Eq.~(\ref{eq:rh_j_identity}).  Hilbert--P\'olya compatibility:
Ihara categorically excluded; spectral compatible at current sample
precision.  The two sides are not two views of the same object;
they are two distinct framework observables computed from the
same Dirac-$S^3$ data, with structural support that is disjoint
in the bipartite-$\chi_{-4}$ sense.

\paragraph{Origin: Paper~24's Coulomb/HO asymmetry.}
The meta-pattern is consistent with---and arguably a consequence
of---the Coulomb/HO asymmetry established in
Paper~24~\cite{loutey_paper24}: the graph is rational/integer on
both bundles, but the projection structures (second-order Riemannian
for the scalar, first-order complex-analytic for the HO holomorphic
sector, Dirac for the spinor) are categorically different, and
their transcendental projections onto physical coordinates do not
interpolate.  Paper~25~\cite{loutey_paper25} sharpens this further:
abelian U(1) Wilson--Hodge structure transfers verbatim from $S^3$
to the Bargmann--Segal $S^5$ graph, but non-abelian SU(3) on $(N,0)$
SU(3) irreps fails (Clebsch--Gordan intertwiners are not group
elements), making non-abelian gauge content Coulomb-specific.  The
universal/Coulomb partition is therefore three-layered, not two:
(i)~spectrum-computing role of $L_0$ (yes $S^3$, no $S^5$);
(ii)~calibration $\pi$ (yes $S^3$, no $S^5$);
(iii)~Wilson physical content (full $S^3$, reduced $S^5$).

\paragraph{Operational consequence.}  For taxonomy construction
(Sec.~\ref{sec:taxonomy}), the consequence is that multi-component
observables should be decomposed and each component classified
separately, rather than the whole observable forced into a single
tier.  This changes nothing about
Sec.~\ref{sec:observable_classification}'s S/P/C classification of
atomic observables (which are fine as they stand at the observable
level), but it does mean that composite structural results---such as
Paper~2's $K$ or Paper~28's two-loop sunset---should be broken into
their components before taxonomic assignment.

% ------------------------------------------------------------------

%% ========================================================================
%% IX. CONCLUSION
%% ========================================================================

\section{Conclusion}
\label{sec:conclusion}

The exchange-constant taxonomy of Sec.~\ref{sec:taxonomy} now
admits three theorems and one observation, alongside the
catalogued constants of Levels~1--5:

\paragraph{Theorem 1 (operator-order discriminant).}
On round~$S^3$, the squared Dirac operator $D^2$ has spectral zeta
$\zeta_{D^2}(s) = 2^{2s-1}[\lambda(2s{-}2) - \lambda(2s)]$ which is
a polynomial in~$\pi^2$ at every integer~$s \ge 2$
(Eq.~\ref{eq:zeta_d2}); the first-order Dirac operator has
spectral zeta containing odd-zeta content
$\zeta_R(2k{+}1) \perp \mathbb{Q}(\pi)$ at odd~$s$
(Eq.~\ref{eq:hurwitz_discriminant}).  Operator order is the primary
transcendental discriminant; bundle type modulates rational
coefficients only.  Loop-order prediction (verified): no odd-zeta
at one loop (Schwinger~\cite{schwinger1948}, $a_e^{(1)} = \alpha/2\pi$);
$\zeta_R(3)$ at two loops (Petermann~\cite{petermann1957};
Sommerfield~\cite{sommerfield1957}).

\paragraph{Theorem 2 (three-axis classification).}
The transcendental content of any spectral functional on round~$S^3$
is determined by three independent axes:
(A1)~operator order (1st vs.\ 2nd) via Theorem~\ref{thm:operator_order};
(A2)~bundle type (scalar vs.\ spinor) modulating coefficients;
(A3)~diagram topology (free vs.\ vertex-restricted) introducing
Dirichlet $L$-values ($G = \beta(2)$, $\beta(4)$) through the
quarter-integer Hurwitz shift.  The classification is $R$-independent:
$D(s,R)/R^s$ is a structural invariant, and Weyl's law preserves the
even/odd-$s$ discriminant in the flat-space limit.

\paragraph{Theorem 3 (three independent spectral homes for $K$).}
The three components of Paper~2's $K = \pi(B + F - \Delta)$ have
structurally independent spectral origins:
$B$ a finite Casimir trace on the scalar Laplacian at $n_{\max}=3$;
$F = D_{n^2}(d_{\max}) = \zeta_R(2)$ an infinite Dirichlet series
of Fock degeneracies at the packing exponent;
$\Delta^{-1} = g_3^{\mathrm{Dirac}} = 40$ a finite Dirac mode
degeneracy on the spinor bundle.  No common generator produces all
three.  The combination rule $K = \alpha^{-1}$
(Observation~\ref{obs:k_rule}) remains a numerical observation
without first-principles derivation.

\paragraph{Observation (structural incommensurability).}
The three theorems share a meta-pattern: framework observables that
look unitary at the surface
(``$K = \pi(B + F - \Delta)$'';
``two-loop QED on $S^3$'';
``zero structure of the GeoVac zeta'')
resolve under structural analysis into objects from different cells
of the operator-order $\times$ bundle $\times$ topology grid.
Observation~\ref{obs:incommensurability} catalogues three independent
instances; no further unification is currently visible.

\paragraph{Verified and open.}
The taxonomy has been validated across the full GeoVac hierarchy
through Level~5; Claim~4 (transcendental content determined by
projection type) has held across all tested cases.  What remains
open: (i)~whether the three-axis classification of
Theorem~\ref{thm:three_axis} predicts the physical two-loop
coefficient (not just the transcendental class); (ii)~the
combination rule $K = \pi(B + F - \Delta)$ of
Observation~\ref{obs:k_rule}, whose component decomposition is now
Theorem~\ref{thm:k_decomposition} but whose additive conspiracy
remains unexplained; (iii)~whether the structural incommensurability
of Observation~\ref{obs:incommensurability} sharpens to a theorem
when further composite observables are decomposed.

\paragraph{Design principle.}
Staying on the graph eliminates exchange constants by eliminating
projections.  The program's effort to algebraicize solvers (Paper~12's
Neumann expansion, Laguerre recurrences, Gaunt replacements,
Paper~19's cross-center $V_{ne}$ via incomplete gamma functions) is,
in retrospect, an effort to eliminate exchange constants.  The algebraic
curve result (Sec.~\ref{sec:algebraic_curve}) shows this has been more
successful than previously recognized.

\paragraph{Note added v2.6.0.}
The He exchange constant census (Track~DI) validates the taxonomy
computationally.  Graph-native FCI achieves 0.19\% at $n_{\max} = 7$
with zero free parameters.  Three-layer structure: (1)~rational
(graph $h_1$ + Slater $V_{ee}$, 98.6\%); (2)~topological (graph
off-diagonal $T_\pm$, calibration $\kappa$, dominant); (3)~transcendental
(e-e cusp, embedding, ${\sim}0.14\%$ floor).  FCI basis invariance
confirmed: the floor is embedding content (cusp), not basis mismatch.

\paragraph{Note added v2.26.1 (Seeley--DeWitt cancellation).}
On unit $S^{3}$, the Seeley--DeWitt coefficient ratio
$a_{k}^{D^{2}} / a_{k}^{\Delta_{\mathrm{LB}}} = 4 = \dim(S)$
at every order $k$, where $S$ denotes the spinor bundle.  This ratio
is exact: the asymptotic mode density of the Camporesi--Higuchi Dirac
spectrum is $g_{n}^{\mathrm{Dirac}} \sim 4n^{2}$, four times the
scalar density $n^{2}$, and the SD expansion depends only on this
asymptotic density.  As a consequence, the perturbative
Connes--Chamseddine supertrace
$\mathrm{Str}[f(D^{2}/\Lambda^{2})] \equiv
 \mathrm{Tr}_{S}[f] - \tfrac{1}{4}\mathrm{Tr}_{D}[f]$
vanishes identically in the asymptotic expansion.  The three
ingredients of $K/\pi = B + F - \Delta$ are all invisible to the
SD expansion: $B$ is a finite Casimir trace, $F$ is an arithmetic
Dirichlet series, and $\Delta^{-1} = 40$ is the Euler--Maclaurin
upper boundary term of the Dirac mode-count sum at $n_{\mathrm{CH}}
= 3$.  The sign on $\Delta$ in the combination rule $K = \pi(B + F -
\Delta)$ is the standard $(-1)^{F}$ grading: $B$ and $F$ are
scalar-sector (bosonic) quantities that enter with $(+)$, while
$\Delta$ is a spinor-sector (fermionic) quantity that enters with
$(-)$.  $K/\pi$ lives entirely in the non-perturbative sector of the
spectral-action supertrace on~$S^{3}$.

\begin{acknowledgments}
Computational implementation and documentation drafting were
performed with the assistance of large language models (Anthropic
Claude) under the author's scientific direction.  All results
were validated against analytical benchmarks.
\end{acknowledgments}


%% ========================================================================
%% BIBLIOGRAPHY
%% ========================================================================

\begin{thebibliography}{30}

\bibitem{loutey_paper0}
J.~Loutey,
``Angular Momentum as Information Geometry: Isotropic Packing and the Universal Angular Cross-Section,''
GeoVac Paper~0 (2026).

\bibitem{loutey_paper1}
J.~Loutey,
``The Geometric Atom: Atomic Structure as a Packing Problem,''
GeoVac Paper~1 (2026).

\bibitem{loutey_paper7}
J.~Loutey,
``The Dimensionless Vacuum: Recovering the Schr\"odinger Equation
from Scale-Invariant Graph Topology,''
GeoVac Paper~7 (2026).

\bibitem{loutey_paper11}
J.~Loutey,
``The Molecular Fock Projection: Diatomic Molecules as Prolate
Spheroidal Lattices,''
GeoVac Paper~11 (2026).

\bibitem{loutey_paper13}
J.~Loutey,
``The Hyperspherical Lattice: Two-Electron Atoms as Coupled
Channel Graphs,''
GeoVac Paper~13 (2026).

\bibitem{loutey_fci}
J.~Loutey,
``Full configuration interaction on a sparse graph Hamiltonian:
Multi-electron atoms via relational lattice indexing,''
GeoVac FCI Paper (2026).

\bibitem{loutey_paper2}
J.~Loutey,
``The Fine Structure Constant as a Spectral Coincidence on the Hopf
Fibration,''
GeoVac Paper~2 (2026), \texttt{papers/group5\_qed\_gauge/}.

\bibitem{loutey_paper22}
J.~Loutey,
``Potential-Independent Angular Sparsity for Qubit Hamiltonians in
Spherical Fermion Systems,''
GeoVac Paper~22 (2026).

\bibitem{loutey_paper24}
J.~Loutey,
``The Bargmann--Segal Lattice: $\pi$-Free Discretization of the
Harmonic Oscillator on $S^5$,''
GeoVac Paper~24 (2026).

\bibitem{loutey_paper25}
J.~Loutey,
``The Hopf Graph as a Lattice Gauge Structure: A Framework Observation,''
GeoVac Paper~25 (2026).

\bibitem{loutey_paper27}
J.~Loutey,
``Entropy as a Projection Artifact: One-Body Operators are
Entanglement-Inert on Sparse Lattices,''
GeoVac Paper~27 (2026).

\bibitem{loutey_paper28}
J.~Loutey,
``QED on $S^3$: Spectral Geometry of Perturbative Transcendentals,''
GeoVac Paper~28 (2026).

\bibitem{loutey_paper29}
J.~Loutey,
``The GeoVac Hopf Graphs Are Ramanujan: Ihara Zeta, Graph-RH, and
a Scope Boundary on Selberg-on-Hydrogen,''
GeoVac Paper~29 (2026).

\bibitem{loutey_paper30}
J.~Loutey,
``SU(2) Wilson Lattice Gauge Structure on the GeoVac Hopf Graph:
The Non-Abelian Sibling of Paper~25,''
GeoVac Paper~30 (2026).

\bibitem{loutey_paper32}
J.~Loutey,
``The GeoVac Spectral Triple: A Synthesis Paper Locking the Framework
into the Marcolli--van~Suijlekom Gauge-Network Lineage,''
GeoVac Paper~32 (2026).

\bibitem{loutey_paper31}
J.~Loutey,
``The Universal/Coulomb Partition of the GeoVac Spectral Triple,''
GeoVac Paper~31 (2026).

\bibitem{loutey_paper55}
J.~Loutey,
``Periods of GeoVac: a cyclotomic mixed-Tate classification of the
master Mellin engine,''
GeoVac Paper~55 (2026).

\bibitem{krajewski1998}
T.~Krajewski,
``Classification of Finite Spectral Triples,''
J. Geom. Phys. \textbf{28}, 1--30 (1998),
\href{https://arxiv.org/abs/hep-th/9701081}{arXiv:hep-th/9701081}.

\bibitem{paschke_sitarz2000}
M.~Paschke and A.~Sitarz,
``Discrete Spectral Triples and their Symmetries,''
J. Math. Phys. \textbf{39}, 6191--6205 (1998),
\href{https://arxiv.org/abs/q-alg/9612029}{arXiv:q-alg/9612029}.

\bibitem{loutey_track_j}
J.~Loutey,
``Track J: Algebraic Level 2 Radial Solver for $m \neq 0$ States,''
GeoVac supplementary material (2026).

\bibitem{loutey_track_g}
J.~Loutey,
``Track G: Perturbation Series for $\mu(R)$,''
GeoVac supplementary material (2026).

\bibitem{loutey_track_m}
J.~Loutey,
``Track M: Prolate CI $l_{\max}$ Convergence Study,''
GeoVac supplementary material (2026).

\bibitem{loutey_track_p1}
J.~Loutey,
``Track P1: Algebraic Curve for $\mu(R)$---Characteristic Polynomial
of the Level~3 Angular Hamiltonian,''
GeoVac supplementary material (2026).

\bibitem{loutey_track_s}
J.~Loutey,
``Track S: Level~4 Algebraic Curve---Structural Negative Result,''
GeoVac supplementary material (2026).

\bibitem{drake2006}
G.~W.~F.~Drake,
``High precision calculations for helium,''
in \textit{Springer Handbook of Atomic, Molecular, and Optical Physics},
edited by G.~W.~F.~Drake (Springer, New York, 2006), pp.~199--219.

\bibitem{dlmf_laguerre}
NIST Digital Library of Mathematical Functions,
\S18.9.13, \url{https://dlmf.nist.gov/18.9.13}.

\bibitem{weyl1911}
H.~Weyl,
``\"Uber die asymptotische Verteilung der Eigenwerte,''
Nachr.\ K\"onigl.\ Ges.\ Wiss.\ G\"ottingen, Math.-phys.\ Kl.\
(1911), pp.\ 110--117.

\bibitem{selberg1956}
A.~Selberg,
``Harmonic analysis and discontinuous groups in weakly symmetric
Riemannian spaces with applications to Dirichlet series,''
J.\ Indian Math.\ Soc.\ \textbf{20}, 47--87 (1956).

\bibitem{berger2003}
M.~Berger,
\textit{A Panoramic View of Riemannian Geometry}
(Springer, Berlin, 2003).

\bibitem{lichnerowicz1963}
A.~Lichnerowicz,
\textit{Spineurs harmoniques},
C.~R.\ Acad.\ Sci.\ Paris \textbf{257}, 7--9 (1963).

\bibitem{camporesi_higuchi1996}
R.~Camporesi and A.~Higuchi,
On the eigenfunctions of the Dirac operator on spheres and
real hyperbolic spaces,
\textit{J.\ Geom.\ Phys.}\ \textbf{20}, 1--18 (1996).

\bibitem{peter_weyl1927}
F.~Peter and H.~Weyl,
``Die Vollst\"andigkeit der primitiven Darstellungen einer
geschlossenen kontinuierlichen Gruppe,''
\textit{Math.\ Ann.}\ \textbf{97}, 737--755 (1927).

\bibitem{polchinski1998}
J.~Polchinski,
\textit{String Theory}, Vol.~2: Superstring Theory and Beyond
(Cambridge University Press, Cambridge, 1998).

\bibitem{goncharov1999}
A.~B.~Goncharov,
``Volumes of hyperbolic manifolds and mixed Tate motives,''
\textit{J. Amer. Math. Soc.} \textbf{12}, 569--618 (1999).

\bibitem{zagier1994}
D.~Zagier,
``Values of zeta functions and their applications,''
in \textit{First European Congress of Mathematics}, Vol.~II
(Birkh\"auser, Basel, 1994), pp.~497--512.

\bibitem{schwinger1948}
J.~Schwinger,
``On quantum-electrodynamics and the magnetic moment of the electron,''
\textit{Phys.\ Rev.}\ \textbf{73}, 416--417 (1948).

\bibitem{petermann1957}
A.~Petermann,
``Fourth order magnetic moment of the electron,''
\textit{Helv.\ Phys.\ Acta} \textbf{30}, 407--408 (1957).

\bibitem{sommerfield1957}
C.~M.~Sommerfield,
``Magnetic dipole moment of the electron,''
\textit{Phys.\ Rev.}\ \textbf{107}, 328--329 (1957).

\bibitem{adamchik2003}
V.~S.~Adamchik,
``A certain series associated with Catalan's constant,''
\textit{Z.\ Anal.\ Anwend.}\ \textbf{22}, 817--826 (2003).

\bibitem{seeley1967}
R.~T.~Seeley,
``Complex powers of an elliptic operator,''
\textit{Proc.\ Sympos.\ Pure Math.}\ \textbf{10}, 288--307 (1967).

\bibitem{minakshisundaram1949}
S.~Minakshisundaram and {\AA}.~Pleijel,
``Some properties of the eigenfunctions of the Laplace-operator
on Riemannian manifolds,''
\textit{Canad.\ J.\ Math.}\ \textbf{1}, 242--256 (1949).

\bibitem{gilkey1975}
P.~B.~Gilkey,
``The spectral geometry of a Riemannian manifold,''
\textit{J.\ Differential Geom.}\ \textbf{10}, 601--618 (1975).

\bibitem{latremoliere2018}
F.~Latr\'emoli\`ere, ``The Dual Modular Propinquity and Completeness,''
arXiv:1811.04534;
\textit{J.\ Noncommut.\ Geom.}\ \textbf{15}, 347 (2021).

\bibitem{fathizadeh_marcolli2016}
F.~Fathizadeh and M.~Marcolli,
``Periods and motives in the spectral action of Robertson--Walker spacetimes,''
arXiv:1611.01815 (2016).

\bibitem{deligne2010}
P.~Deligne,
``Le groupe fondamental unipotent motivique de
$\mathbb{G}_m - \mu_N$, pour $N = 2, 3, 4, 6$ ou $8$,''
\textit{Publ.~Math.~IHES} \textbf{112}, 101--141 (2010).
% NB: arXiv:math/0302267 is the separate Deligne--Goncharov paper, not this one.

\bibitem{glanois2015}
C.~Glanois,
``Motivic unipotent fundamental groupoid of
$\mathbb{G}_m \setminus \mu_N$ for $N = 2, 3, 4, 6, 8$ and Galois descents,''
arXiv:1411.4947 (v2, 2 Sep 2015);
published version \textit{J.~Number Theory} \textbf{160}, 334--384 (2016).

\bibitem{rejzner2016}
K.~Rejzner,
``Renormalization and periods in perturbative Algebraic Quantum Field Theory,''
arXiv:1603.02748 (2016).

\bibitem{connes_marcolli2004}
A.~Connes and M.~Marcolli,
``Renormalization and motivic Galois theory,''
arXiv:math/0409306 (2004).

\bibitem{karamata1962}
J.~Karamata,
``\"Uber die Hardy--Littlewoodschen Umkehrungen des Abelschen
Stetigkeitssatzes,''
\textit{Math.\ Z.}\ \textbf{32}, 319--320 (1930);
reprint in \textit{Selected Papers}, Mathematical Institute,
Belgrade, 1962, \S V.3.

\bibitem{korevaar2004}
J.~Korevaar,
\textit{Tauberian Theory: A Century of Developments},
Grundlehren der mathematischen Wissenschaften \textbf{329},
Springer-Verlag, Berlin (2004).

\bibitem{marcolli_vs2014}
M.~Marcolli and W.~D.~van~Suijlekom,
``Gauge networks in noncommutative geometry,''
\textit{J.~Geom.\ Phys.}\ \textbf{75}, 71--91 (2014);
arXiv:1301.3480.

\bibitem{perez_sanchez2024}
C.~I.~Perez-Sanchez,
``Bratteli networks and the Spectral Action on quivers,''
arXiv:2401.03705 (2024).

\bibitem{perez_sanchez2025}
C.~I.~Perez-Sanchez,
``Comment on `Gauge networks in noncommutative geometry','',
arXiv:2508.17338 (2025).

\bibitem{brown2012}
F.~Brown,
``Mixed Tate motives over $\mathbb{Z}$,''
\textit{Ann.\ of Math.}\ (2)~\textbf{175}, 949--976 (2012).

\end{thebibliography}

\end{document}
